% GENERATED FILE -- do not edit by hand.
% Regenerate with: make index   (tools/gen_question_index.py)
\chapter{Question Index}
\label{chap:question_index}

\begin{tldr}
Every one of the 301 interview questions in Volume I, indexed by chapter, by round, and by frequency. This appendix is generated from the questions themselves, so it cannot drift from the chapters. Work the \freqhigh{} list first: those are the questions that come up in almost every loop.
\end{tldr}

\section{At a Glance}
\begin{table}[H]\centering
\caption{Question distribution}
\begin{tabular}{lr}
\toprule
\textbf{Category} & \textbf{Count} \\
\midrule
\quad \textbf{Total} & \textbf{301} \\
\quad L5 questions & 98 \\
\quad L6 questions & 160 \\
\quad L7 questions & 43 \\
\quad Breadth round & 75 \\
\quad Depth round & 153 \\
\quad Coding round & 16 \\
\quad Design round & 50 \\
\quad Behavioral round & 7 \\
\quad \freqhigh{} asked constantly & 151 \\
\quad \freqmed{} common & 124 \\
\quad \freqlow{} occasional & 26 \\
\bottomrule
\end{tabular}
\end{table}

\section{Start Here: The \freqhigh{} Questions}
These come up in almost every loop at this level. If you have one evening, these are the evening.
\begin{enumerate}[leftmargin=*]
\item \textbf{Ch.~1} \quad L5 \quad \textit{Conceptual} \quad Explain SVD and give three applications in machine learning
\item \textbf{Ch.~1} \quad L5 \quad \textit{First Principles} \quad Why does the chain rule matter for backpropagation? Walk through a 3-layer example
\item \textbf{Ch.~1} \quad L5 \quad \textit{Conceptual} \quad What is KL divergence? Why is it not a true distance metric? When do you use it in practice?
\item \textbf{Ch.~1} \quad L5 \quad \textit{Estimation} \quad Estimate the memory for storing a 50K $\times$ 768 embedding matrix in FP32, FP16, and INT8
\item \textbf{Ch.~1} \quad L6 \quad \textit{First Principles} \quad Derive the gradient of softmax cross-entropy loss with respect to the logits
\item \textbf{Ch.~2} \quad L5 \quad \textit{Conceptual} \quad Explain the bias-variance tradeoff. Does it apply to deep learning?
\item \textbf{Ch.~2} \quad L5 \quad \textit{Conceptual} \quad Explain double descent. Why does test error decrease again in the over-parameterized regime?
\item \textbf{Ch.~2} \quad L6 \quad \textit{Trade-off} \quad Your 10B parameter model is clearly overparameterized for your task. Should you reduce it?
\item \textbf{Ch.~2} \quad L6 \quad \textit{Trade-off} \quad Your team has a fixed compute budget. How do you decide model size vs. training data?
\item \textbf{Ch.~2} \quad L6 \quad \textit{Estimation} \quad Estimate the compute-optimal model size for a training budget of $10^{22}$ FLOPs using Chinchilla scaling laws
\item \textbf{Ch.~3} \quad L5 \quad \textit{Conceptual} \quad When do you use softmax cross-entropy vs.\ sigmoid BCE?
\item \textbf{Ch.~3} \quad L5 \quad \textit{Trade-off} \quad Cross-entropy vs.\ focal loss vs.\ class-balanced loss for imbalanced classification---what is your decision framework?
\item \textbf{Ch.~3} \quad L6 \quad \textit{Debugging} \quad Your SFT loss changes when you change batch composition or gradient-accumulation steps---same data, same model, same effective batch size. Why?
\item \textbf{Ch.~3} \quad L5 \quad \textit{Conceptual} \quad Why does SimCLR need large batch sizes, and what are the alternatives?
\item \textbf{Ch.~3} \quad L6 \quad \textit{Architecture Design} \quad You're building a visual search system for an e-commerce platform with 10M products. What loss function would you use for training the embedding model?
\item \textbf{Ch.~3} \quad L5 \quad \textit{Architecture Design} \quad You need to train an embedding model for semantic search across 100M documents. Which loss function and why?
\item \textbf{Ch.~3} \quad L6 \quad \textit{Trade-off} \quad Pointwise vs.\ pairwise vs.\ listwise ranking losses---when do you use each?
\item \textbf{Ch.~3} \quad L6 \quad \textit{Architecture Design} \quad Design the loss function for a retrieval system where you have click data (implicit feedback) but no explicit relevance labels
\item \textbf{Ch.~3} \quad L5 \quad \textit{Trade-off} \quad Your classification model gets 95\% accuracy but the loss curve shows the model is still decreasing slowly. Should you keep training?
\item \textbf{Ch.~3} \quad L5 \quad \textit{Architecture Design} \quad You're building a product recommendation system with implicit feedback (clicks). What loss would you use and why?
\item \textbf{Ch.~4} \quad L5 \quad \textit{Conceptual} \quad Explain the dying ReLU problem. How do modern activations solve it?
\item \textbf{Ch.~4} \quad L5 \quad \textit{Trade-off} \quad Why GELU over ReLU? Why SwiGLU over GELU?
\item \textbf{Ch.~4} \quad L5 \quad \textit{Debugging} \quad Training loss is NaN after a few hundred steps. Which activation/normalization issues could cause this?
\item \textbf{Ch.~4} \quad L6 \quad \textit{Debugging} \quad Your 70B-parameter run in BF16 starts loss-spiking at 100B tokens. Walk me through the activation- and normalization-level causes and mitigations
\item \textbf{Ch.~5} \quad L6 \quad \textit{Trade-off} \quad Compare MHA, GQA, MQA, and MLA for KV-cache efficiency. Why does MLA need a decoupled RoPE key?
\item \textbf{Ch.~5} \quad L5 \quad \textit{Conceptual} \quad Explain Flash Attention. Why is it faster even though it does the same computation?
\item \textbf{Ch.~5} \quad L6 \quad \textit{Estimation} \quad Estimate the memory required to train a 7B parameter model with batch size 1, sequence length 2048
\item \textbf{Ch.~5} \quad L5 \quad \textit{Conceptual} \quad How does RoPE encode position information? Why is it preferred over learned positional embeddings?
\item \textbf{Ch.~5} \quad L5 \quad \textit{Trade-off} \quad Compare MHA, GQA, and MQA. When would you choose each?
\item \textbf{Ch.~5} \quad L5 \quad \textit{First Principles} \quad Why does attention scale by $1/\sqrt{d_k}$? What happens if you don't?
\item \textbf{Ch.~5} \quad L6 \quad \textit{Debugging} \quad Your large BF16 pretraining run shows intermittent loss spikes. Instrumentation shows attention logits growing into the hundreds over training. Diagnose and fix
\item \textbf{Ch.~5} \quad L6 \quad \textit{Estimation} \quad Calculate the FLOPs for a single forward pass through a 7B parameter Transformer
\item \textbf{Ch.~5} \quad L6 \quad \textit{Estimation} \quad Calculate the KV cache memory for serving LLaMA 70B at 128K context length
\item \textbf{Ch.~5} \quad L5 \quad \textit{Conceptual} \quad What changed between GPT-2 (2019) and LLaMA (2023) architecturally? Why each change?
\item \textbf{Ch.~6} \quad L6 \quad \textit{Estimation} \quad You are putting 100M item embeddings behind a vector index. Which embedding-side choices do you make---dimension, precision, refresh---and how would you know if one of them silently cost you recall?
\item \textbf{Ch.~6} \quad L5 \quad \textit{Conceptual} \quad How does subword tokenization (BPE) interact with embedding quality?
\item \textbf{Ch.~6} \quad L5 \quad \textit{Debugging} \quad How would you debug poor embedding quality?
\item \textbf{Ch.~7} \quad L5 \quad \textit{Debugging} \quad Your two-tower model retrieves documents that share keywords with the query but aren't actually relevant. How do you fix this?
\item \textbf{Ch.~7} \quad L5 \quad \textit{Trade-off} \quad Siamese vs Triplet vs Two-Tower---decision framework for a new retrieval project
\item \textbf{Ch.~7} \quad L5 \quad \textit{Trade-off} \quad Cross-encoder vs bi-encoder---when is the quality gap worth the latency cost?
\item \textbf{Ch.~7} \quad L5 \quad \textit{Conceptual} \quad Cosine similarity vs dot product vs L2 distance for embeddings---when each?
\item \textbf{Ch.~8} \quad L6 \quad \textit{Architecture Design} \quad Design an SSL pre-training pipeline for a dataset of 10M unlabeled images and 100K labeled images. Walk through method selection, architecture, training, and evaluation
\item \textbf{Ch.~8} \quad L5 \quad \textit{Trade-off} \quad SimCLR vs BYOL vs MAE---give me a decision framework for choosing between them
\item \textbf{Ch.~8} \quad L5 \quad \textit{Trade-off} \quad BERT vs GPT pre-training objectives---what are the fundamental tradeoffs between masked language modeling and autoregressive language modeling?
\item \textbf{Ch.~8} \quad L6 \quad \textit{First Principles} \quad Next-token prediction vs masked language modeling---why did GPT's autoregressive approach win for generative AI while BERT's bidirectional approach initially dominated NLU benchmarks?
\item \textbf{Ch.~8} \quad L5 \quad \textit{Conceptual} \quad Why do non-contrastive methods like BYOL work without negative samples? Shouldn't they collapse to a trivial solution?
\item \textbf{Ch.~8} \quad L6 \quad \textit{Architecture Design} \quad You have 10M unlabeled images and 10K labeled images in a medical imaging domain. How would you use SSL to improve your classifier, and what challenges do you expect?
\item \textbf{Ch.~9} \quad L5 \quad \textit{Estimation} \quad You're told a model is ``Llama-class, 7B parameters.'' Roughly where do those 7B parameters live?
\item \textbf{Ch.~9} \quad L5 \quad \textit{Trade-off} \quad Encoder-only vs decoder-only vs encoder-decoder---when do you use each, and what drives the decision?
\item \textbf{Ch.~9} \quad L6 \quad \textit{Debugging} \quad Your fine-tuned chat model performs worse than the base model it was tuned from. Walk me through your diagnosis
\item \textbf{Ch.~9} \quad L6 \quad \textit{Debugging} \quad Your fine-tuned LLM generates fluent but factually wrong answers. Diagnose the issue and propose solutions
\item \textbf{Ch.~9} \quad L6 \quad \textit{Conceptual} \quad Why do decoder-only models dominate modern NLP despite encoder-decoder being theoretically more flexible?
\item \textbf{Ch.~9} \quad L6 \quad \textit{First Principles} \quad Walk me through the DPO loss function. How does it eliminate the need for a separate reward model?
\item \textbf{Ch.~9} \quad L6 \quad \textit{Architecture Design} \quad Design the architecture for a customer support ticket classification system with 50 categories, handling multilingual input and evolving category taxonomy
\item \textbf{Ch.~10} \quad L5 \quad \textit{Trade-off} \quad CNN vs ViT for a dataset with only 10K labeled images---how do you decide, and what changes with 10M images?
\item \textbf{Ch.~10} \quad L5 \quad \textit{Conceptual} \quad Explain how DETR works and why it represents a paradigm shift in object detection
\item \textbf{Ch.~11} \quad L5 \quad \textit{First Principles} \quad Explain the reparameterization trick and why it's needed for VAE training. What happens without it?
\item \textbf{Ch.~11} \quad L6 \quad \textit{Trade-off} \quad Diffusion models vs GANs vs VAEs vs autoregressive models---give me a comparison table and decision framework for different generation tasks
\item \textbf{Ch.~11} \quad L5 \quad \textit{Conceptual} \quad What is classifier-free guidance and why does it improve generation quality in diffusion models?
\item \textbf{Ch.~11} \quad L5 \quad \textit{First Principles} \quad Explain the ELBO derivation for VAEs and why the model produces blurry images. What fixes exist?
\item \textbf{Ch.~11} \quad L5 \quad \textit{Conceptual} \quad How does Stable Diffusion work end-to-end? Walk through the architecture from text prompt to generated image
\item \textbf{Ch.~12} \quad L5 \quad \textit{Conceptual} \quad How do you handle the cold start problem for new users and new items?
\item \textbf{Ch.~12} \quad L5 \quad \textit{Trade-off} \quad Two-tower vs.\ interaction-based models for candidate generation---when would you use each?
\item \textbf{Ch.~14} \quad L6 \quad \textit{Architecture Design} \quad Design an efficient architecture for on-device ML (phone, $<$100ms latency, $<$50MB model)
\item \textbf{Ch.~14} \quad L7 \quad \textit{Architecture Design} \quad You have a serving budget of 80\,GB of GPU memory and 20\,ms/token decode latency, with contexts up to 128K. Design the model
\item \textbf{Ch.~14} \quad L6 \quad \textit{Trade-off} \quad INT8 vs INT4 vs FP8 quantization---what is your decision framework?
\item \textbf{Ch.~15} \quad L5 \quad \textit{Debugging} \quad Your model's loss plateaus after initial descent. Walk through your debugging process
\item \textbf{Ch.~15} \quad L6 \quad \textit{Architecture Design} \quad How would you scale training of a 70B parameter model across 256 GPUs?
\item \textbf{Ch.~15} \quad L5 \quad \textit{Trade-off} \quad Adam vs SGD with momentum---when would you pick each?
\item \textbf{Ch.~15} \quad L5 \quad \textit{Debugging} \quad Training loss is NaN after 1000 steps. Walk through your debugging process
\item \textbf{Ch.~15} \quad L7 \quad \textit{Debugging} \quad Your 70B pretraining run diverges at 200B tokens. Give a full differential diagnosis, ordered by prior probability, with the cheapest discriminating experiment for each
\item \textbf{Ch.~15} \quad L7 \quad \textit{Debugging} \quad A node dies every 3 hours on your 16K-GPU run. Design the recovery system
\item \textbf{Ch.~15} \quad L6 \quad \textit{Estimation} \quad How much GPU memory is needed to fine-tune a 13B parameter model with LoRA vs.\ full fine-tuning?
\item \textbf{Ch.~15} \quad L6 \quad \textit{Trade-off} \quad Data parallelism vs.\ tensor parallelism vs.\ pipeline parallelism---what is your decision framework?
\item \textbf{Ch.~16} \quad L5 \quad \textit{Estimation} \quad Estimate the number of trainable parameters for LoRA with rank 16 on a 7B model
\item \textbf{Ch.~16} \quad L6 \quad \textit{Debugging} \quad Your LoRA fine-tuned model performs well on your test set but catastrophically forgets general capabilities. How do you fix it?
\item \textbf{Ch.~16} \quad L6 \quad \textit{Debugging} \quad Your fine-tuned model echoes the prompt before answering---why?
\item \textbf{Ch.~16} \quad L6 \quad \textit{Trade-off} \quad LoRA vs QLoRA vs full fine-tuning vs prompt tuning---give a complete decision framework
\item \textbf{Ch.~16} \quad L6 \quad \textit{Trade-off} \quad You have chosen LoRA to adapt a 7B instruction-tuned model to an internal support-and-policy domain: 20K curated examples, one 80\,GB GPU, and the adapter will be served alongside the general assistant. Pick rank, $\alpha$, and target modules, and defend the choices
\item \textbf{Ch.~16} \quad L5 \quad \textit{Architecture Design} \quad You need to fine-tune a 70B LLM for a customer's specific domain. You have access to a single A100 80GB GPU. What approach would you take?
\item \textbf{Ch.~16} \quad L6 \quad \textit{Debugging} \quad You fine-tuned a model and it performs well on your test set but poorly in production. What happened?
\item \textbf{Ch.~17} \quad L5 \quad \textit{Conceptual} \quad Walk me through the RLHF pipeline for aligning an LLM
\item \textbf{Ch.~17} \quad L6 \quad \textit{Trade-off} \quad PPO vs DPO vs RLHF---complete comparison for LLM alignment
\item \textbf{Ch.~17} \quad L5 \quad \textit{Conceptual} \quad What is reward hacking? Give 3 concrete examples and how to mitigate each
\item \textbf{Ch.~17} \quad L6 \quad \textit{Debugging} \quad Your RLHF-tuned model becomes sycophantic (agrees with everything). Diagnose and fix
\item \textbf{Ch.~17} \quad L5 \quad \textit{First Principles} \quad Why does the KL penalty matter in RLHF? What happens without it?
\item \textbf{Ch.~17} \quad L6 \quad \textit{First Principles} \quad Explain GRPO. Why did it replace PPO's critic for reasoning RL, and what breaks when you use it naively?
\item \textbf{Ch.~17} \quad L7 \quad \textit{Trade-off} \quad When is test-time compute cheaper than a bigger model---and when does it stop working?
\item \textbf{Ch.~18} \quad L6 \quad \textit{Architecture Design} \quad Design a multimodal content understanding system for a social media platform
\item \textbf{Ch.~18} \quad L6 \quad \textit{Trade-off} \quad Adapter-based (LLaVA) vs natively multimodal (Gemini)---tradeoffs
\item \textbf{Ch.~18} \quad L6 \quad \textit{Debugging} \quad Your VLM hallucinates visual details that are not in the image. Why and how to fix?
\item \textbf{Ch.~18} \quad L5 \quad \textit{Conceptual} \quad Compare early, late, and cross-attention fusion for multimodal models. When would you use each?
\item \textbf{Ch.~19} \quad L7 \quad \textit{Architecture Design} \quad Design an LLM serving stack that handles 1,000 queries per second with a p99 latency of 2 seconds for a 70B parameter model
\item \textbf{Ch.~19} \quad L7 \quad \textit{Architecture Design} \quad Your capacity plan assumed 500-token outputs. The reasoning model you are now deploying emits $\sim$20K thinking tokens per query. Redo the plan
\item \textbf{Ch.~19} \quad L5 \quad \textit{Conceptual} \quad Explain speculative decoding. When would you use it and when would you not?
\item \textbf{Ch.~19} \quad L6 \quad \textit{Debugging} \quad KV cache memory is your bottleneck---you are running out of GPU memory and dropping requests. What do you do?
\item \textbf{Ch.~19} \quad L6 \quad \textit{Estimation} \quad How much memory savings does INT4 quantization give for a 70B model? What is the quality trade-off?
\item \textbf{Ch.~19} \quad L6 \quad \textit{Trade-off} \quad Tensor parallelism vs pipeline parallelism for inference---what is your decision framework?
\item \textbf{Ch.~20} \quad L6 \quad \textit{Architecture Design} \quad Design a content filtering pipeline for an LLM chatbot with $<$50ms added latency
\item \textbf{Ch.~20} \quad L5 \quad \textit{Conceptual} \quad Jailbreak attacks: how do they work and how do you defend?
\item \textbf{Ch.~20} \quad L5 \quad \textit{Conceptual} \quad How do you detect and reduce hallucination in a production LLM system?
\item \textbf{Ch.~20} \quad L7 \quad \textit{Architecture Design} \quad Design safety for an agent that can browse the web and execute code
\item \textbf{Ch.~21} \quad L6 \quad \textit{Trade-off} \quad Your 8K-trained model must serve 64K contexts next month. What are your options, and what do they cost?
\item \textbf{Ch.~21} \quad L6 \quad \textit{Estimation} \quad Estimate the cost per query for a RAG system vs a long-context LLM (128K)
\item \textbf{Ch.~21} \quad L5 \quad \textit{Trade-off} \quad Long context window vs RAG---when do you choose each?
\item \textbf{Ch.~22} \quad L6 \quad \textit{Architecture Design} \quad Design ML infrastructure for a search ranking system at 10K QPS
\item \textbf{Ch.~22} \quad L6 \quad \textit{Debugging} \quad Your model's CTR dropped 5\% overnight. Walk through your diagnosis
\item \textbf{Ch.~22} \quad L6 \quad \textit{Architecture Design} \quad How would you set up an A/B test for a new recommendation model?
\item \textbf{Ch.~22} \quad L5 \quad \textit{Conceptual} \quad How do you handle training-serving skew?
\item \textbf{Ch.~22} \quad L6 \quad \textit{Debugging} \quad Your model's online metrics diverge from offline metrics after 2 weeks in production. What is happening?
\item \textbf{Ch.~22} \quad L5 \quad \textit{Trade-off} \quad Shadow deployment vs A/B test vs multi-armed bandit---when do you use each?
\item \textbf{Ch.~23} \quad L6 \quad \textit{Estimation} \quad Model B is 1.5 points better than model A on MMLU---do you ship it?
\item \textbf{Ch.~23} \quad L7 \quad \textit{Architecture Design} \quad Design the eval stack that decides whether tomorrow's checkpoint ships
\item \textbf{Ch.~23} \quad L5 \quad \textit{Conceptual} \quad Your model has 95\% accuracy but stakeholders are unhappy. What is going on?
\item \textbf{Ch.~23} \quad L5 \quad \textit{Conceptual} \quad How do you evaluate a search ranking system?
\item \textbf{Ch.~23} \quad L5 \quad \textit{Trade-off} \quad Precision is 0.95 but recall is 0.30. What do you do?
\item \textbf{Ch.~23} \quad L5 \quad \textit{Conceptual} \quad How do you know if a model improvement is statistically significant?
\item \textbf{Ch.~23} \quad L6 \quad \textit{Debugging} \quad Your model's AUC is 0.95 but calibration is poor (ECE = 0.15). Why does this matter and how do you fix it?
\item \textbf{Ch.~23} \quad L5 \quad \textit{Trade-off} \quad NDCG vs MAP vs MRR for ranking---when do you use each?
\item \textbf{Ch.~23} \quad L6 \quad \textit{Debugging} \quad Offline metrics went up but online metrics went down. Give 5 possible reasons
\item \textbf{Ch.~23} \quad L5 \quad \textit{Conceptual} \quad What is wrong with accuracy as a metric? When is it actually fine?
\item \textbf{Ch.~24} \quad L6 \quad \textit{Architecture Design} \quad Design a content recommendation system for a news app
\item \textbf{Ch.~24} \quad L6 \quad \textit{Architecture Design} \quad Design the ML architecture for a content moderation system
\item \textbf{Ch.~24} \quad L6 \quad \textit{Trade-off} \quad Your model needs to run in $<$10ms. Current latency is 50ms. Walk through the optimization decision tree
\item \textbf{Ch.~24} \quad L6 \quad \textit{Trade-off} \quad You have 3 months and a team of 3 ML engineers. Should you fine-tune an LLM or train a custom model?
\item \textbf{Ch.~24} \quad L5 \quad \textit{Trade-off} \quad Would you use BERT or GPT for sentiment analysis?
\item \textbf{Ch.~24} \quad L5 \quad \textit{Trade-off} \quad When should you use gradient boosting vs.\ neural networks for tabular data?
\item \textbf{Ch.~25} \quad L5 \quad \textit{Conceptual} \quad Walk me through what happens between ``we have a Common Crawl snapshot'' and the first training step
\item \textbf{Ch.~25} \quad L7 \quad \textit{System Design} \quad You own data for the next frontier pretraining run. Design the pipeline and the ablation program
\item \textbf{Ch.~25} \quad L6 \quad \textit{Trade-off} \quad How would you decide the code:web:math mixture for a pretraining run?
\item \textbf{Ch.~25} \quad L7 \quad \textit{Trade-off} \quad Your data lead proposes raising the synthetic fraction of the next run from 10\% to 40\%. How do you evaluate the proposal?
\item \textbf{Ch.~26} \quad L6 \quad \textit{Estimation} \quad I give you an op---say, a LayerNorm over a $[16{,}384 \times 8{,}192]$ BF16 tensor, or a $4096^3$ GEMM. Is it compute-bound or bandwidth-bound on an H100? Show your method
\item \textbf{Ch.~26} \quad L6 \quad \textit{Debugging} \quad Your training job shows the GPU at 20\% utilization. Diagnose it
\item \textbf{Ch.~26} \quad L7 \quad \textit{Estimation} \quad Estimate the maximum tokens/second for a 70B dense model on 8$\times$H100---for training, and for inference
\item \textbf{Ch.~27} \quad L5 \quad \textit{Conceptual} \quad What makes an LLM system an ``agent''? Walk me through the agent loop and where the engineering difficulty actually lives
\item \textbf{Ch.~27} \quad L6 \quad \textit{Trade-off} \quad When should you NOT build an agent? Your PM wants ``an agent'' for a document-processing product---how do you decide?
\item \textbf{Ch.~27} \quad L6 \quad \textit{System Design} \quad Design a code-review agent for your organization's monorepo
\item \textbf{Ch.~27} \quad L7 \quad \textit{Debugging} \quad Your agent succeeds on 60\% of an internal task suite. Leadership wants 95\%. Take it there
\item \textbf{Ch.~28} \quad L5 \quad \textit{Coding} \quad Implement multi-head self-attention with a causal mask in numpy
\item \textbf{Ch.~28} \quad L5 \quad \textit{Coding} \quad Implement one AdamW update step in numpy
\item \textbf{Ch.~28} \quad L5 \quad \textit{Coding} \quad Write a minimal training loop with gradient clipping in PyTorch
\item \textbf{Ch.~28} \quad L5 \quad \textit{Coding} \quad Implement BPE---train the merges on a corpus, then encode a word
\item \textbf{Ch.~28} \quad L5 \quad \textit{Coding} \quad Implement sampling from logits with temperature, top-$k$, and top-$p$
\item \textbf{Ch.~28} \quad L6 \quad \textit{Coding} \quad Implement greedy decoding with a KV cache
\item \textbf{Ch.~28} \quad L5 \quad \textit{Coding} \quad Implement NDCG@k
\item \textbf{Ch.~28} \quad L5 \quad \textit{Coding} \quad Implement logistic regression with SGD from scratch, deriving the gradient
\item \textbf{Ch.~28} \quad L5 \quad \textit{Coding} \quad Implement k-means with k-means++ initialization
\item \textbf{Ch.~29} \quad L6 \quad \textit{Behavioral} \quad Walk me through your most impactful project. (Then 40 minutes of drilling.)
\item \textbf{Ch.~29} \quad L6 \quad \textit{Behavioral} \quad A peer team's launch is hurting your metric. Their leadership is celebrating it. What do you do?
\item \textbf{Ch.~29} \quad L6 \quad \textit{Behavioral} \quad Describe committing to a technical decision you argued against. Would you do it again?
\item \textbf{Ch.~29} \quad L6 \quad \textit{System Design} \quad I'm going to give you a deliberately vague prompt---``design memory for our assistant.'' Run the first ten minutes
\end{enumerate}

\section{By Chapter}
\subsection*{Chapter 1: Mathematical Foundations}
\begin{tabularx}{\textwidth}{@{}l l l X@{}}
\toprule
\textbf{Lvl} & \textbf{Freq} & \textbf{Type} & \textbf{Question} \\
\midrule
L5 & \freqmed{} & Estimation & Estimate the FLOPs for multiplying two matrices of size $(1024, 4096)$ and $(4096, 1024)$ \\
L5 & \freqmed{} & Conceptual & What are eigenvalues and why do they matter for training neural networks? \\
L5 & \freqhigh{} & Conceptual & Explain SVD and give three applications in machine learning \\
L5 & \freqhigh{} & First Principles & Why does the chain rule matter for backpropagation? Walk through a 3-layer example \\
L6 & \freqmed{} & Debugging & Your gradient norms are 1000$\times$ larger in early layers than later layers. What is happening mathematically? \\
L5 & \freqhigh{} & Conceptual & What is KL divergence? Why is it not a true distance metric? When do you use it in practice? \\
L6 & \freqmed{} & Conceptual & KL divergence is not symmetric. When does this matter in practice? \\
L5 & \freqhigh{} & Estimation & Estimate the memory for storing a 50K $\times$ 768 embedding matrix in FP32, FP16, and INT8 \\
L6 & \freqhigh{} & First Principles & Derive the gradient of softmax cross-entropy loss with respect to the logits \\
L7 & \freqlow{} & First Principles & What is the rank of the attention matrix and why does this matter for efficiency? \\
\bottomrule
\end{tabularx}

\subsection*{Chapter 2: Learning Theory Essentials}
\begin{tabularx}{\textwidth}{@{}l l l X@{}}
\toprule
\textbf{Lvl} & \textbf{Freq} & \textbf{Type} & \textbf{Question} \\
\midrule
L5 & \freqhigh{} & Conceptual & Explain the bias-variance tradeoff. Does it apply to deep learning? \\
L5 & \freqhigh{} & Conceptual & Explain double descent. Why does test error decrease again in the over-parameterized regime? \\
L6 & \freqhigh{} & Trade-off & Your 10B parameter model is clearly overparameterized for your task. Should you reduce it? \\
L6 & \freqhigh{} & Trade-off & Your team has a fixed compute budget. How do you decide model size vs. training data? \\
L6 & \freqhigh{} & Estimation & Estimate the compute-optimal model size for a training budget of $10^{22}$ FLOPs using Chinchilla scaling laws \\
L7 & \freqmed{} & Architecture Design & You get 1\% of the target training budget to run scaling experiments that must predict the final run's loss. Design the study \\
L7 & \freqmed{} & Conceptual & Neural scaling laws: what scales and what does not? When do they break? \\
L6 & \freqmed{} & Conceptual & What is grokking and what does it tell us about generalization? \\
L6 & \freqmed{} & Conceptual & Are emergent abilities in LLMs real, or an artifact of evaluation? \\
L6 & \freqmed{} & Trade-off & Training loss is still decreasing but validation loss plateaued. You have $10\times$ more unlabeled data available. What do you do? \\
L7 & \freqmed{} & First Principles & Why do large language models generalize despite having more parameters than training examples? \\
\bottomrule
\end{tabularx}

\subsection*{Chapter 3: Loss Functions: The Interview Playbook}
\begin{tabularx}{\textwidth}{@{}l l l X@{}}
\toprule
\textbf{Lvl} & \textbf{Freq} & \textbf{Type} & \textbf{Question} \\
\midrule
L5 & \freqhigh{} & Conceptual & When do you use softmax cross-entropy vs.\ sigmoid BCE? \\
L6 & \freqmed{} & Debugging & Your model's training loss is decreasing but validation loss curves show the model is increasingly miscalibrated. What loss function issue could cause this? \\
L5 & \freqhigh{} & Trade-off & Cross-entropy vs.\ focal loss vs.\ class-balanced loss for imbalanced classification---what is your decision framework? \\
L6 & \freqhigh{} & Debugging & Your SFT loss changes when you change batch composition or gradient-accumulation steps---same data, same model, same effective batch size. Why? \\
L6 & \freqmed{} & Estimation & Estimate the memory and FLOP cost of the cross-entropy layer for a 128K-vocabulary model training on a 1M-token batch. How do you reduce it? \\
L5 & \freqhigh{} & Conceptual & Why does SimCLR need large batch sizes, and what are the alternatives? \\
L6 & \freqhigh{} & Architecture Design & You're building a visual search system for an e-commerce platform with 10M products. What loss function would you use for training the embedding model? \\
L6 & \freqlow{} & Debugging & Your contrastive learning model's loss suddenly jumps mid-training and doesn't recover. Diagnose \\
L5 & \freqhigh{} & Architecture Design & You need to train an embedding model for semantic search across 100M documents. Which loss function and why? \\
L6 & \freqhigh{} & Trade-off & Pointwise vs.\ pairwise vs.\ listwise ranking losses---when do you use each? \\
L6 & \freqhigh{} & Architecture Design & Design the loss function for a retrieval system where you have click data (implicit feedback) but no explicit relevance labels \\
L5 & \freqhigh{} & Trade-off & Your classification model gets 95\% accuracy but the loss curve shows the model is still decreasing slowly. Should you keep training? \\
L5 & \freqhigh{} & Architecture Design & You're building a product recommendation system with implicit feedback (clicks). What loss would you use and why? \\
L7 & \freqmed{} & Architecture Design & You're building a multi-task model that does classification, regression, and ranking simultaneously. Design the loss function \\
\bottomrule
\end{tabularx}

\subsection*{Chapter 4: Activation Functions and Normalization}
\begin{tabularx}{\textwidth}{@{}l l l X@{}}
\toprule
\textbf{Lvl} & \textbf{Freq} & \textbf{Type} & \textbf{Question} \\
\midrule
L5 & \freqhigh{} & Conceptual & Explain the dying ReLU problem. How do modern activations solve it? \\
L5 & \freqhigh{} & Trade-off & Why GELU over ReLU? Why SwiGLU over GELU? \\
L6 & \freqmed{} & Debugging & Your model shows different behavior at training vs inference. What normalization issue could cause this? \\
L6 & \freqmed{} & First Principles & Why don't Transformers use BatchNorm? Derive the reasoning from first principles \\
L5 & \freqmed{} & Trade-off & RMSNorm vs LayerNorm---when would you choose each? \\
L6 & \freqmed{} & Trade-off & Pre-norm vs. post-norm---derive the gradient-flow difference and explain what modern hybrids fix \\
L5 & \freqhigh{} & Debugging & Training loss is NaN after a few hundred steps. Which activation/normalization issues could cause this? \\
L6 & \freqhigh{} & Debugging & Your 70B-parameter run in BF16 starts loss-spiking at 100B tokens. Walk me through the activation- and normalization-level causes and mitigations \\
\bottomrule
\end{tabularx}

\subsection*{Chapter 5: Attention Mechanisms and Transformers}
\begin{tabularx}{\textwidth}{@{}l l l X@{}}
\toprule
\textbf{Lvl} & \textbf{Freq} & \textbf{Type} & \textbf{Question} \\
\midrule
L6 & \freqhigh{} & Trade-off & Compare MHA, GQA, MQA, and MLA for KV-cache efficiency. Why does MLA need a decoupled RoPE key? \\
L5 & \freqhigh{} & Conceptual & Explain Flash Attention. Why is it faster even though it does the same computation? \\
L6 & \freqhigh{} & Estimation & Estimate the memory required to train a 7B parameter model with batch size 1, sequence length 2048 \\
L5 & \freqhigh{} & Conceptual & How does RoPE encode position information? Why is it preferred over learned positional embeddings? \\
L5 & \freqhigh{} & Trade-off & Compare MHA, GQA, and MQA. When would you choose each? \\
L5 & \freqhigh{} & First Principles & Why does attention scale by $1/\sqrt{d_k}$? What happens if you don't? \\
L6 & \freqhigh{} & Debugging & Your large BF16 pretraining run shows intermittent loss spikes. Instrumentation shows attention logits growing into the hundreds over training. Diagnose and fix \\
L6 & \freqhigh{} & Estimation & Calculate the FLOPs for a single forward pass through a 7B parameter Transformer \\
L6 & \freqhigh{} & Estimation & Calculate the KV cache memory for serving LLaMA 70B at 128K context length \\
L5 & \freqmed{} & First Principles & Why is self-attention permutation equivariant? Why does this matter for Transformers? \\
L5 & \freqmed{} & Trade-off & Pre-norm vs Post-norm Transformers --- which would you choose and why? \\
L6 & \freqmed{} & Debugging & Your Transformer model shows degrading perplexity at 32K context despite being trained on 8K. Diagnose and fix \\
L7 & \freqlow{} & Architecture Design & Design the attention mechanism for a model that must process 1 million token documents \\
L5 & \freqhigh{} & Conceptual & What changed between GPT-2 (2019) and LLaMA (2023) architecturally? Why each change? \\
L6 & \freqmed{} & Trade-off & Compare Transformer vs SSM (Mamba) architectures. When would you choose each? \\
L7 & \freqlow{} & Debugging & Your model's attention patterns show all heads attending to the same positions (typically position 0 or EOS). What's happening? \\
L7 & \freqlow{} & Estimation & Estimate the ratio of attention FLOPs to FFN FLOPs at different sequence lengths. When does attention become the bottleneck? \\
L7 & \freqlow{} & Architecture Design & You need to modify a standard Transformer for real-time streaming applications (e.g., live transcription). What changes? \\
L6 & \freqmed{} & Trade-off & Flash Attention vs Linear Attention --- when would you use each? \\
\bottomrule
\end{tabularx}

\subsection*{Chapter 6: Embeddings and Representation Learning}
\begin{tabularx}{\textwidth}{@{}l l l X@{}}
\toprule
\textbf{Lvl} & \textbf{Freq} & \textbf{Type} & \textbf{Question} \\
\midrule
L6 & \freqmed{} & Debugging & Your embedding space shows that unrelated items cluster together. Diagnose and fix \\
L6 & \freqhigh{} & Estimation & You are putting 100M item embeddings behind a vector index. Which embedding-side choices do you make---dimension, precision, refresh---and how would you know if one of them silently cost you recall? \\
L5 & \freqmed{} & Estimation & Estimate the memory for an embedding table with 50K vocab, 768 dimensions in FP16 vs INT8 \\
L6 & \freqmed{} & Conceptual & What is embedding collapse and how do you detect/prevent it? \\
L5 & \freqhigh{} & Conceptual & How does subword tokenization (BPE) interact with embedding quality? \\
L5 & \freqhigh{} & Debugging & How would you debug poor embedding quality? \\
\bottomrule
\end{tabularx}

\subsection*{Chapter 7: Similarity and Metric Learning Architectures}
\begin{tabularx}{\textwidth}{@{}l l l X@{}}
\toprule
\textbf{Lvl} & \textbf{Freq} & \textbf{Type} & \textbf{Question} \\
\midrule
L5 & \freqhigh{} & Debugging & Your two-tower model retrieves documents that share keywords with the query but aren't actually relevant. How do you fix this? \\
L5 & \freqhigh{} & Trade-off & Siamese vs Triplet vs Two-Tower---decision framework for a new retrieval project \\
L6 & \freqmed{} & Debugging & Your two-tower model's recall@10 is good but precision@10 is poor. What's wrong? \\
L5 & \freqhigh{} & Trade-off & Cross-encoder vs bi-encoder---when is the quality gap worth the latency cost? \\
L6 & \freqmed{} & Trade-off & Hard negative mining strategies---how to choose and when each fails \\
L5 & \freqhigh{} & Conceptual & Cosine similarity vs dot product vs L2 distance for embeddings---when each? \\
L7 & \freqlow{} & Debugging & Your contrastive model produces good embeddings for frequent items but poor for rare items. Why? \\
L6 & \freqmed{} & Estimation & Estimate the QPS of a retrieval system using HNSW on 100M documents. Then add a cross-encoder reranker over the top 100---what changes? \\
\bottomrule
\end{tabularx}

\subsection*{Chapter 8: Self-Supervised Learning}
\begin{tabularx}{\textwidth}{@{}l l l X@{}}
\toprule
\textbf{Lvl} & \textbf{Freq} & \textbf{Type} & \textbf{Question} \\
\midrule
L5 & \freqmed{} & Conceptual & What makes a good augmentation policy for contrastive learning, and how does the choice of augmentations define the learned representation? \\
L6 & \freqhigh{} & Architecture Design & Design an SSL pre-training pipeline for a dataset of 10M unlabeled images and 100K labeled images. Walk through method selection, architecture, training, and evaluation \\
L5 & \freqhigh{} & Trade-off & SimCLR vs BYOL vs MAE---give me a decision framework for choosing between them \\
L6 & \freqmed{} & Debugging & Your SSL model's downstream performance plateaus despite increasing pre-training data from 1M to 10M images. What could be going wrong? \\
L5 & \freqhigh{} & Trade-off & BERT vs GPT pre-training objectives---what are the fundamental tradeoffs between masked language modeling and autoregressive language modeling? \\
L6 & \freqmed{} & Estimation & Estimate the compute cost of pre-training a BERT-base model on 16GB of text data \\
L6 & \freqhigh{} & First Principles & Next-token prediction vs masked language modeling---why did GPT's autoregressive approach win for generative AI while BERT's bidirectional approach initially dominated NLU benchmarks? \\
L5 & \freqhigh{} & Conceptual & Why do non-contrastive methods like BYOL work without negative samples? Shouldn't they collapse to a trivial solution? \\
L6 & \freqhigh{} & Architecture Design & You have 10M unlabeled images and 10K labeled images in a medical imaging domain. How would you use SSL to improve your classifier, and what challenges do you expect? \\
L6 & \freqmed{} & First Principles & Explain the InfoNCE loss and its connection to mutual information. Why does temperature matter? \\
\bottomrule
\end{tabularx}

\subsection*{Chapter 9: NLP Architectures}
\begin{tabularx}{\textwidth}{@{}l l l X@{}}
\toprule
\textbf{Lvl} & \textbf{Freq} & \textbf{Type} & \textbf{Question} \\
\midrule
L5 & \freqhigh{} & Estimation & You're told a model is ``Llama-class, 7B parameters.'' Roughly where do those 7B parameters live? \\
L5 & \freqhigh{} & Trade-off & Encoder-only vs decoder-only vs encoder-decoder---when do you use each, and what drives the decision? \\
L6 & \freqhigh{} & Debugging & Your fine-tuned chat model performs worse than the base model it was tuned from. Walk me through your diagnosis \\
L6 & \freqhigh{} & Debugging & Your fine-tuned LLM generates fluent but factually wrong answers. Diagnose the issue and propose solutions \\
L6 & \freqmed{} & Estimation & Estimate the inference latency for generating 100 tokens with a 7B parameter model on a single A100 GPU \\
L7 & \freqlow{} & First Principles & How does in-context learning work mechanistically? Why can large language models learn from examples in the prompt without any gradient updates? \\
L5 & \freqmed{} & Trade-off & BPE vs SentencePiece vs character-level tokenization---what are the tradeoffs and how does tokenizer choice affect model performance? \\
L6 & \freqhigh{} & Conceptual & Why do decoder-only models dominate modern NLP despite encoder-decoder being theoretically more flexible? \\
L7 & \freqmed{} & Architecture Design & Design the architecture for a multilingual document understanding system that handles 50+ languages, mixed-language documents, and both structured (tables, forms) and unstructured text \\
L6 & \freqhigh{} & First Principles & Walk me through the DPO loss function. How does it eliminate the need for a separate reward model? \\
L6 & \freqhigh{} & Architecture Design & Design the architecture for a customer support ticket classification system with 50 categories, handling multilingual input and evolving category taxonomy \\
L6 & \freqmed{} & Conceptual & Compare RoPE vs learned positional embeddings vs sinusoidal encodings. Why did RoPE become the dominant choice for modern LLMs? \\
\bottomrule
\end{tabularx}

\subsection*{Chapter 10: Vision Architectures}
\begin{tabularx}{\textwidth}{@{}l l l X@{}}
\toprule
\textbf{Lvl} & \textbf{Freq} & \textbf{Type} & \textbf{Question} \\
\midrule
L5 & \freqhigh{} & Trade-off & CNN vs ViT for a dataset with only 10K labeled images---how do you decide, and what changes with 10M images? \\
L6 & \freqmed{} & Estimation & Estimate the FLOPs for a ResNet-50 forward pass on a $224 \times 224$ image. How does this compare to ViT-B/16? \\
L6 & \freqmed{} & Architecture Design & Design a vision backbone for real-time object detection on an edge device (e.g., mobile phone or drone with $<$10W power budget) \\
L5 & \freqmed{} & Conceptual & What is the effective receptive field of a CNN and why does it matter for architecture design? \\
L6 & \freqmed{} & Debugging & Your ViT model achieves 92\% accuracy on ImageNet but only 65\% on your domain-specific dataset despite fine-tuning. Diagnose the problem \\
L6 & \freqmed{} & Architecture Design & Design a visual search system that handles 100M product images with sub-200ms latency. Walk through the architecture \\
L5 & \freqhigh{} & Conceptual & Explain how DETR works and why it represents a paradigm shift in object detection \\
L6 & \freqmed{} & Trade-off & Compare the ResNet skip connection to DenseNet dense connections to Swin Transformer's shifted windows. What problem does each solve? \\
\bottomrule
\end{tabularx}

\subsection*{Chapter 11: Generative Models}
\begin{tabularx}{\textwidth}{@{}l l l X@{}}
\toprule
\textbf{Lvl} & \textbf{Freq} & \textbf{Type} & \textbf{Question} \\
\midrule
L5 & \freqhigh{} & First Principles & Explain the reparameterization trick and why it's needed for VAE training. What happens without it? \\
L6 & \freqhigh{} & Trade-off & Diffusion models vs GANs vs VAEs vs autoregressive models---give me a comparison table and decision framework for different generation tasks \\
L6 & \freqmed{} & Debugging & Your diffusion model generates high-quality images overall but has mode collapse for certain categories (e.g., generates dogs well but all cats look the same). Diagnose and fix \\
L6 & \freqmed{} & Architecture Design & Design an image generation pipeline for a product photography application (e.g., placing products in realistic scenes) \\
L5 & \freqhigh{} & Conceptual & What is classifier-free guidance and why does it improve generation quality in diffusion models? \\
L5 & \freqhigh{} & First Principles & Explain the ELBO derivation for VAEs and why the model produces blurry images. What fixes exist? \\
L7 & \freqlow{} & Estimation & Estimate the inference cost (FLOPs and latency) for generating a $512 \times 512$ image with Stable Diffusion on an A100 GPU \\
L7 & \freqlow{} & Conceptual & Flow matching vs denoising diffusion---what changed and why is it considered the next evolution? \\
L5 & \freqhigh{} & Conceptual & How does Stable Diffusion work end-to-end? Walk through the architecture from text prompt to generated image \\
L6 & \freqmed{} & First Principles & Explain the WGAN loss and why Wasserstein distance solves GAN training instability \\
\bottomrule
\end{tabularx}

\subsection*{Chapter 12: Recommendation Systems}
\begin{tabularx}{\textwidth}{@{}l l l X@{}}
\toprule
\textbf{Lvl} & \textbf{Freq} & \textbf{Type} & \textbf{Question} \\
\midrule
L5 & \freqhigh{} & Conceptual & How do you handle the cold start problem for new users and new items? \\
L5 & \freqhigh{} & Trade-off & Two-tower vs.\ interaction-based models for candidate generation---when would you use each? \\
\bottomrule
\end{tabularx}

\subsection*{Chapter 13: Graph Neural Networks}
\begin{tabularx}{\textwidth}{@{}l l l X@{}}
\toprule
\textbf{Lvl} & \textbf{Freq} & \textbf{Type} & \textbf{Question} \\
\midrule
L6 & \freqmed{} & Architecture Design & Design a GNN-based system for real-time fraud detection on a financial transaction graph \\
L6 & \freqmed{} & Trade-off & When should you NOT use a GNN? Give examples where simpler approaches win \\
L5 & \freqmed{} & Debugging & Your GNN's performance degrades as you add more layers. What's happening and how do you fix it? \\
L5 & \freqmed{} & Trade-off & GCN vs GraphSAGE vs GAT---when would you use each? \\
L6 & \freqlow{} & Estimation & Estimate the memory requirements for running a 3-layer GCN on a graph with 10M nodes \\
L6 & \freqmed{} & Trade-off & GCN vs GraphSAGE vs GAT---which would you choose for a production recommendation system at scale? \\
L6 & \freqlow{} & First Principles & What are the theoretical expressivity limitations of message-passing GNNs, and when do they matter in practice? \\
L5 & \freqmed{} & Conceptual & How do you handle heterogeneous graphs (multiple node types, multiple edge types) in a GNN? \\
\bottomrule
\end{tabularx}

\subsection*{Chapter 14: Efficient Architectures}
\begin{tabularx}{\textwidth}{@{}l l l X@{}}
\toprule
\textbf{Lvl} & \textbf{Freq} & \textbf{Type} & \textbf{Question} \\
\midrule
L6 & \freqhigh{} & Architecture Design & Design an efficient architecture for on-device ML (phone, $<$100ms latency, $<$50MB model) \\
L6 & \freqmed{} & Trade-off & When would you choose MoE over dense scaling, and vice versa? \\
L6 & \freqmed{} & Trade-off & You're building a system that processes 100K-token documents. Compare transformer vs Mamba for this use case \\
L6 & \freqmed{} & Estimation & Estimate the compute savings of a MoE model with 8 experts, top-2 routing, vs a dense model of the same quality \\
L7 & \freqhigh{} & Architecture Design & You have a serving budget of 80\,GB of GPU memory and 20\,ms/token decode latency, with contexts up to 128K. Design the model \\
L6 & \freqhigh{} & Trade-off & INT8 vs INT4 vs FP8 quantization---what is your decision framework? \\
L6 & \freqmed{} & Debugging & Your quantized model performs well on benchmarks but fails on specific domains (code, math). Why? \\
L6 & \freqmed{} & Architecture Design & Knowledge distillation: design a pipeline to compress a 70B teacher to a 7B student \\
L6 & \freqmed{} & Debugging & MoE routing collapse: what is it, how do you detect it, and how do you fix it? \\
L6 & \freqmed{} & Trade-off & Structured vs unstructured pruning---which actually gives real speedup? \\
L7 & \freqlow{} & Architecture Design & Design an efficient architecture for real-time video understanding on edge devices \\
\bottomrule
\end{tabularx}

\subsection*{Chapter 15: Training Optimization}
\begin{tabularx}{\textwidth}{@{}l l l X@{}}
\toprule
\textbf{Lvl} & \textbf{Freq} & \textbf{Type} & \textbf{Question} \\
\midrule
L5 & \freqhigh{} & Debugging & Your model's loss plateaus after initial descent. Walk through your debugging process \\
L6 & \freqhigh{} & Architecture Design & How would you scale training of a 70B parameter model across 256 GPUs? \\
L5 & \freqhigh{} & Trade-off & Adam vs SGD with momentum---when would you pick each? \\
L5 & \freqmed{} & Conceptual & Explain gradient accumulation and when you would use it \\
L6 & \freqmed{} & Debugging & Your model's loss spikes every N steps then recovers. Diagnose \\
L5 & \freqhigh{} & Debugging & Training loss is NaN after 1000 steps. Walk through your debugging process \\
L7 & \freqhigh{} & Debugging & Your 70B pretraining run diverges at 200B tokens. Give a full differential diagnosis, ordered by prior probability, with the cheapest discriminating experiment for each \\
L6 & \freqmed{} & Debugging & Your distributed training job is 60\% as fast as expected with 8 GPUs. What is wrong? \\
L7 & \freqhigh{} & Debugging & A node dies every 3 hours on your 16K-GPU run. Design the recovery system \\
L6 & \freqhigh{} & Estimation & How much GPU memory is needed to fine-tune a 13B parameter model with LoRA vs.\ full fine-tuning? \\
L7 & \freqlow{} & Estimation & Calculate the training throughput (tokens/sec) for a 7B model on 8$\times$A100s \\
L5 & \freqmed{} & Trade-off & Cosine decay vs.\ constant LR with warmup---when would you use each? \\
L6 & \freqhigh{} & Trade-off & Data parallelism vs.\ tensor parallelism vs.\ pipeline parallelism---what is your decision framework? \\
L6 & \freqmed{} & Trade-off & FP16 vs.\ BF16 vs.\ FP8 for training---when would you use each? \\
\bottomrule
\end{tabularx}

\subsection*{Chapter 16: Transfer Learning and Parameter-Efficient Fine-Tuning}
\begin{tabularx}{\textwidth}{@{}l l l X@{}}
\toprule
\textbf{Lvl} & \textbf{Freq} & \textbf{Type} & \textbf{Question} \\
\midrule
L5 & \freqhigh{} & Estimation & Estimate the number of trainable parameters for LoRA with rank 16 on a 7B model \\
L6 & \freqhigh{} & Debugging & Your LoRA fine-tuned model performs well on your test set but catastrophically forgets general capabilities. How do you fix it? \\
L6 & \freqhigh{} & Debugging & Your fine-tuned model echoes the prompt before answering---why? \\
L6 & \freqhigh{} & Trade-off & LoRA vs QLoRA vs full fine-tuning vs prompt tuning---give a complete decision framework \\
L6 & \freqhigh{} & Trade-off & You have chosen LoRA to adapt a 7B instruction-tuned model to an internal support-and-policy domain: 20K curated examples, one 80\,GB GPU, and the adapter will be served alongside the general assistant. Pick rank, $\alpha$, and target modules, and defend the choices \\
L6 & \freqmed{} & Architecture Design & Design a fine-tuning pipeline for adapting a foundation model to a new domain with 10K examples \\
L5 & \freqmed{} & Conceptual & What is negative transfer and how do you detect and prevent it? \\
L6 & \freqmed{} & First Principles & Why does LoRA work? What is the low-rank assumption about fine-tuning? \\
L5 & \freqmed{} & Trade-off & Adapter layers vs LoRA vs prefix tuning---what are the architectural differences and when would you use each? \\
L5 & \freqhigh{} & Architecture Design & You need to fine-tune a 70B LLM for a customer's specific domain. You have access to a single A100 80GB GPU. What approach would you take? \\
L6 & \freqhigh{} & Debugging & You fine-tuned a model and it performs well on your test set but poorly in production. What happened? \\
L6 & \freqmed{} & Architecture Design & You have one base model and 100 customers, each with their own LoRA adapter. Design the serving system \\
\bottomrule
\end{tabularx}

\subsection*{Chapter 17: Reinforcement Learning and RLHF}
\begin{tabularx}{\textwidth}{@{}l l l X@{}}
\toprule
\textbf{Lvl} & \textbf{Freq} & \textbf{Type} & \textbf{Question} \\
\midrule
L5 & \freqhigh{} & Conceptual & Walk me through the RLHF pipeline for aligning an LLM \\
L6 & \freqhigh{} & Trade-off & PPO vs DPO vs RLHF---complete comparison for LLM alignment \\
L5 & \freqhigh{} & Conceptual & What is reward hacking? Give 3 concrete examples and how to mitigate each \\
L6 & \freqhigh{} & Debugging & Your RLHF-tuned model becomes sycophantic (agrees with everything). Diagnose and fix \\
L6 & \freqmed{} & Estimation & Estimate the memory required to run RLHF (4 models: policy, reference, reward model, critic) \\
L5 & \freqhigh{} & First Principles & Why does the KL penalty matter in RLHF? What happens without it? \\
L6 & \freqmed{} & Architecture Design & Design a reward model training pipeline. What data do you need? \\
L6 & \freqmed{} & Trade-off & Constitutional AI vs RLHF vs DPO---when would you use each? \\
L7 & \freqlow{} & Debugging & Your RLHF training shows the reward score plateauing while KL divergence keeps climbing. What is happening? \\
L6 & \freqmed{} & Trade-off & Process reward models vs outcome reward models---what are the tradeoffs and when does each win? \\
L6 & \freqhigh{} & First Principles & Explain GRPO. Why did it replace PPO's critic for reasoning RL, and what breaks when you use it naively? \\
L6 & \freqmed{} & Debugging & Your GRPO run's mean reward climbs but responses grow unboundedly long and benchmark scores stall---diagnose \\
L7 & \freqhigh{} & Trade-off & When is test-time compute cheaper than a bigger model---and when does it stop working? \\
\bottomrule
\end{tabularx}

\subsection*{Chapter 18: Multimodal Learning}
\begin{tabularx}{\textwidth}{@{}l l l X@{}}
\toprule
\textbf{Lvl} & \textbf{Freq} & \textbf{Type} & \textbf{Question} \\
\midrule
L6 & \freqmed{} & Trade-off & Cascade vs native speech-to-speech for a voice assistant---decide and defend \\
L6 & \freqhigh{} & Architecture Design & Design a multimodal content understanding system for a social media platform \\
L6 & \freqhigh{} & Trade-off & Adapter-based (LLaVA) vs natively multimodal (Gemini)---tradeoffs \\
L5 & \freqmed{} & First Principles & How does CLIP's contrastive loss create a shared embedding space? \\
L5 & \freqmed{} & Conceptual & CLIP vs SigLIP---what changed and why? \\
L7 & \freqlow{} & Estimation & Estimate the compute cost of training a CLIP model on 400M image-text pairs \\
L6 & \freqhigh{} & Debugging & Your VLM hallucinates visual details that are not in the image. Why and how to fix? \\
L6 & \freqmed{} & Trade-off & Visual tokens consume 576 tokens of context. How do you reduce this cost? \\
L5 & \freqhigh{} & Conceptual & Compare early, late, and cross-attention fusion for multimodal models. When would you use each? \\
L6 & \freqmed{} & Architecture Design & Design a multimodal search system that lets users search a product catalog using text or image queries \\
L6 & \freqlow{} & Architecture Design & How would you add audio understanding to an existing vision-language model? \\
\bottomrule
\end{tabularx}

\subsection*{Chapter 19: Inference Optimization and LLM Serving}
\begin{tabularx}{\textwidth}{@{}l l l X@{}}
\toprule
\textbf{Lvl} & \textbf{Freq} & \textbf{Type} & \textbf{Question} \\
\midrule
L7 & \freqhigh{} & Architecture Design & Design an LLM serving stack that handles 1,000 queries per second with a p99 latency of 2 seconds for a 70B parameter model \\
L7 & \freqhigh{} & Architecture Design & Your capacity plan assumed 500-token outputs. The reasoning model you are now deploying emits $\sim$20K thinking tokens per query. Redo the plan \\
L5 & \freqhigh{} & Conceptual & Explain speculative decoding. When would you use it and when would you not? \\
L6 & \freqhigh{} & Debugging & KV cache memory is your bottleneck---you are running out of GPU memory and dropping requests. What do you do? \\
L7 & \freqmed{} & Estimation & Calculate the throughput (tokens/sec) of a 70B model on 8$\times$H100s with tensor parallelism \\
L6 & \freqhigh{} & Estimation & How much memory savings does INT4 quantization give for a 70B model? What is the quality trade-off? \\
L7 & \freqlow{} & Estimation & Your LLM service needs to handle 500 concurrent users with 2-second time-to-first-token. Size the infrastructure \\
L6 & \freqmed{} & Trade-off & GPTQ vs AWQ vs GGUF quantization---when would you use each? \\
L6 & \freqhigh{} & Trade-off & Tensor parallelism vs pipeline parallelism for inference---what is your decision framework? \\
L7 & \freqlow{} & Trade-off & Prefill-optimized vs decode-optimized serving---when should you disaggregate them? \\
L6 & \freqmed{} & Debugging & Your LLM serving system's p99 latency is 10$\times$ the p50. Diagnose \\
L6 & \freqmed{} & Debugging & After quantizing your model to INT4, certain types of prompts give garbage output while others are fine. Why? \\
L6 & \freqlow{} & Debugging & Your speculative decoding setup is slower than standard decoding. What went wrong? \\
\bottomrule
\end{tabularx}

\subsection*{Chapter 20: Safety, Alignment, and Interpretability}
\begin{tabularx}{\textwidth}{@{}l l l X@{}}
\toprule
\textbf{Lvl} & \textbf{Freq} & \textbf{Type} & \textbf{Question} \\
\midrule
L6 & \freqhigh{} & Architecture Design & Design a content filtering pipeline for an LLM chatbot with $<$50ms added latency \\
L5 & \freqhigh{} & Conceptual & Jailbreak attacks: how do they work and how do you defend? \\
L6 & \freqmed{} & Conceptual & What is mechanistic interpretability? Give a concrete example of a discovered circuit \\
L7 & \freqmed{} & Trade-off & Your LLM refuses too many benign queries. How do you reduce over-refusal without increasing safety risk? \\
L6 & \freqmed{} & Architecture Design & Hallucination detection: design a system that flags when an LLM generates unsupported claims \\
L5 & \freqhigh{} & Conceptual & How do you detect and reduce hallucination in a production LLM system? \\
L6 & \freqlow{} & Architecture Design & How do you set up a red teaming program for an LLM product? \\
L7 & \freqlow{} & Conceptual & How can you use representation engineering (steering vectors) to control model behavior without retraining? \\
L7 & \freqhigh{} & Architecture Design & Design safety for an agent that can browse the web and execute code \\
L6 & \freqmed{} & Debugging & A jailbreak goes viral against your production model---walk me through the first 24 hours \\
\bottomrule
\end{tabularx}

\subsection*{Chapter 21: Long Context and RAG Systems}
\begin{tabularx}{\textwidth}{@{}l l l X@{}}
\toprule
\textbf{Lvl} & \textbf{Freq} & \textbf{Type} & \textbf{Question} \\
\midrule
L6 & \freqhigh{} & Trade-off & Your 8K-trained model must serve 64K contexts next month. What are your options, and what do they cost? \\
L6 & \freqmed{} & Debugging & Your model's quality degrades beyond 32K context. Walk me through the diagnosis \\
L6 & \freqhigh{} & Estimation & Estimate the cost per query for a RAG system vs a long-context LLM (128K) \\
L5 & \freqhigh{} & Trade-off & Long context window vs RAG---when do you choose each? \\
\bottomrule
\end{tabularx}

\subsection*{Chapter 22: Production ML Systems}
\begin{tabularx}{\textwidth}{@{}l l l X@{}}
\toprule
\textbf{Lvl} & \textbf{Freq} & \textbf{Type} & \textbf{Question} \\
\midrule
L6 & \freqhigh{} & Architecture Design & Design ML infrastructure for a search ranking system at 10K QPS \\
L6 & \freqhigh{} & Debugging & Your model's CTR dropped 5\% overnight. Walk through your diagnosis \\
L6 & \freqhigh{} & Architecture Design & How would you set up an A/B test for a new recommendation model? \\
L5 & \freqhigh{} & Conceptual & How do you handle training-serving skew? \\
L6 & \freqhigh{} & Debugging & Your model's online metrics diverge from offline metrics after 2 weeks in production. What is happening? \\
L7 & \freqmed{} & Estimation & Estimate the infrastructure cost for serving a recommendation model at 100K QPS \\
L6 & \freqmed{} & Architecture Design & Design the monitoring system for an ML model in production \\
L5 & \freqhigh{} & Trade-off & Shadow deployment vs A/B test vs multi-armed bandit---when do you use each? \\
L5 & \freqmed{} & Conceptual & How do you detect and handle concept drift vs data drift? \\
L6 & \freqmed{} & Architecture Design & Feature store: batch vs streaming features---describe the architecture and tradeoffs \\
\bottomrule
\end{tabularx}

\subsection*{Chapter 23: Evaluation Metrics}
\begin{tabularx}{\textwidth}{@{}l l l X@{}}
\toprule
\textbf{Lvl} & \textbf{Freq} & \textbf{Type} & \textbf{Question} \\
\midrule
L6 & \freqhigh{} & Estimation & Model B is 1.5 points better than model A on MMLU---do you ship it? \\
L7 & \freqhigh{} & Architecture Design & Design the eval stack that decides whether tomorrow's checkpoint ships \\
L5 & \freqhigh{} & Conceptual & Your model has 95\% accuracy but stakeholders are unhappy. What is going on? \\
L5 & \freqhigh{} & Conceptual & How do you evaluate a search ranking system? \\
L5 & \freqhigh{} & Trade-off & Precision is 0.95 but recall is 0.30. What do you do? \\
L5 & \freqhigh{} & Conceptual & How do you know if a model improvement is statistically significant? \\
L6 & \freqhigh{} & Debugging & Your model's AUC is 0.95 but calibration is poor (ECE = 0.15). Why does this matter and how do you fix it? \\
L5 & \freqhigh{} & Trade-off & NDCG vs MAP vs MRR for ranking---when do you use each? \\
L6 & \freqmed{} & Estimation & Estimate the sample size needed for an A/B test to detect a 1\% improvement in CTR \\
L7 & \freqmed{} & Architecture Design & Design an evaluation framework for an LLM-based chatbot \\
L6 & \freqhigh{} & Debugging & Offline metrics went up but online metrics went down. Give 5 possible reasons \\
L5 & \freqhigh{} & Conceptual & What is wrong with accuracy as a metric? When is it actually fine? \\
\bottomrule
\end{tabularx}

\subsection*{Chapter 24: Decision Frameworks: When to Use What}
\begin{tabularx}{\textwidth}{@{}l l l X@{}}
\toprule
\textbf{Lvl} & \textbf{Freq} & \textbf{Type} & \textbf{Question} \\
\midrule
L6 & \freqhigh{} & Architecture Design & Design a content recommendation system for a news app \\
L6 & \freqhigh{} & Architecture Design & Design the ML architecture for a content moderation system \\
L6 & \freqhigh{} & Trade-off & Your model needs to run in $<$10ms. Current latency is 50ms. Walk through the optimization decision tree \\
L6 & \freqhigh{} & Trade-off & You have 3 months and a team of 3 ML engineers. Should you fine-tune an LLM or train a custom model? \\
L5 & \freqmed{} & Trade-off & GPU vs TPU vs CPU for inference---give a decision framework for different workloads \\
L7 & \freqmed{} & Trade-off & Build vs buy vs open-source---give a decision framework for ML infrastructure \\
L5 & \freqhigh{} & Trade-off & Would you use BERT or GPT for sentiment analysis? \\
L5 & \freqhigh{} & Trade-off & When should you use gradient boosting vs.\ neural networks for tabular data? \\
\bottomrule
\end{tabularx}

\subsection*{Chapter 25: Pretraining at Scale}
\begin{tabularx}{\textwidth}{@{}l l l X@{}}
\toprule
\textbf{Lvl} & \textbf{Freq} & \textbf{Type} & \textbf{Question} \\
\midrule
L5 & \freqhigh{} & Conceptual & Walk me through what happens between ``we have a Common Crawl snapshot'' and the first training step \\
L7 & \freqhigh{} & System Design & You own data for the next frontier pretraining run. Design the pipeline and the ablation program \\
L6 & \freqhigh{} & Trade-off & How would you decide the code:web:math mixture for a pretraining run? \\
L6 & \freqmed{} & Debugging & After a data-pipeline change, training loss improved---but downstream evals got worse. Diagnose \\
L5 & \freqmed{} & Conceptual & Why does deduplication improve language models, and what goes wrong if you deduplicate too aggressively? \\
L6 & \freqmed{} & System Design & Design the benchmark decontamination strategy for a pretraining run. What does n-gram overlap miss, and what do you do about it? \\
L6 & \freqmed{} & Conceptual & What is mid-training (annealing), why does everyone do it now, and what is the catch? \\
L7 & \freqhigh{} & Trade-off & Your data lead proposes raising the synthetic fraction of the next run from 10\% to 40\%. How do you evaluate the proposal? \\
L6 & \freqmed{} & Trade-off & You are choosing between a 32K, 128K, and 256K vocabulary for a new model family. What changes, and when would you retrain vs.\ reuse a tokenizer? \\
L6 & \freqmed{} & System Design & Design the in-flight evaluation for a three-month pretraining run. When do in-loop evals mislead? \\
L6 & \freqlow{} & Debugging & At 40\% through a flagship run, held-out code loss jumps and stays elevated. Walk me through the response \\
\bottomrule
\end{tabularx}

\subsection*{Chapter 26: GPU Performance Fundamentals}
\begin{tabularx}{\textwidth}{@{}l l l X@{}}
\toprule
\textbf{Lvl} & \textbf{Freq} & \textbf{Type} & \textbf{Question} \\
\midrule
L5 & \freqmed{} & Conceptual & Walk me through the GPU memory hierarchy. Why does it, rather than FLOPs, dominate ML performance thinking? \\
L5 & \freqmed{} & First Principles & Why do tensor cores only accelerate matrix multiplication? What does that imply for how you design models and kernels? \\
L6 & \freqhigh{} & Estimation & I give you an op---say, a LayerNorm over a $[16{,}384 \times 8{,}192]$ BF16 tensor, or a $4096^3$ GEMM. Is it compute-bound or bandwidth-bound on an H100? Show your method \\
L6 & \freqhigh{} & Debugging & Your training job shows the GPU at 20\% utilization. Diagnose it \\
L6 & \freqmed{} & First Principles & Kernel fusion doesn't reduce FLOPs. Why does it make models faster---and when does it stop helping? Use FlashAttention as your example \\
L6 & \freqmed{} & Estimation & How long does it take to all-reduce the gradients of a 70B-parameter model across 8 GPUs? Derive it \\
L6 & \freqlow{} & Debugging & A teammate changed the hidden size from 4096 to 4100 ``to add a few features,'' and training throughput dropped 35\%. What happened? \\
L6 & \freqlow{} & Trade-off & What problem do CUDA graphs solve, and what constraints do they impose? Where do they matter most in an LLM system? \\
L7 & \freqhigh{} & Estimation & Estimate the maximum tokens/second for a 70B dense model on 8$\times$H100---for training, and for inference \\
L7 & \freqmed{} & Debugging & Your 512-GPU pretraining run reports 25\% MFU. Walk me through how you find the missing performance \\
L7 & \freqmed{} & Architecture Design & You have 8 nodes of 8$\times$H100. Map tensor, pipeline, data, and expert parallelism onto this cluster's fabric, justifying the mapping from link characteristics---then tell me what changes on a GB200 NVL72 \\
\bottomrule
\end{tabularx}

\subsection*{Chapter 27: Agents and Tool Use}
\begin{tabularx}{\textwidth}{@{}l l l X@{}}
\toprule
\textbf{Lvl} & \textbf{Freq} & \textbf{Type} & \textbf{Question} \\
\midrule
L5 & \freqhigh{} & Conceptual & What makes an LLM system an ``agent''? Walk me through the agent loop and where the engineering difficulty actually lives \\
L6 & \freqhigh{} & Trade-off & When should you NOT build an agent? Your PM wants ``an agent'' for a document-processing product---how do you decide? \\
L6 & \freqhigh{} & System Design & Design a code-review agent for your organization's monorepo \\
L7 & \freqhigh{} & Debugging & Your agent succeeds on 60\% of an internal task suite. Leadership wants 95\%. Take it there \\
L6 & \freqmed{} & System Design & Design the evaluation for a customer-support agent that can read accounts, process refunds, and escalate to humans \\
L5 & \freqmed{} & Conceptual & Explain pass@$k$ vs.\ pass\textasciicircum{}$k$. Your agent team reports 92\% pass@5---what do you tell the VP who wants to ship? \\
L6 & \freqmed{} & Trade-off & One agent with 30 tools, or an orchestrator with specialized sub-agents? Argue both sides with numbers \\
L6 & \freqmed{} & Conceptual & Why do computer-use agents lag API-based agents so badly, and when would you deploy one anyway? \\
L7 & \freqmed{} & System Design & Your agent's tasks run for hours and blow past the context window. Design its memory and state management \\
L6 & \freqmed{} & Trade-off & Your agent product works but loses money: \$4 of inference per task against \$1.50 of revenue. Fix the economics without wrecking the success rate \\
L5 & \freqmed{} & Conceptual & What does MCP actually standardize, and what does it change---and not change---about how you architect an agent? \\
\bottomrule
\end{tabularx}

\subsection*{Chapter 28: ML Coding Rounds}
\begin{tabularx}{\textwidth}{@{}l l l X@{}}
\toprule
\textbf{Lvl} & \textbf{Freq} & \textbf{Type} & \textbf{Question} \\
\midrule
L5 & \freqhigh{} & Coding & Implement multi-head self-attention with a causal mask in numpy \\
L5 & \freqhigh{} & Coding & Implement one AdamW update step in numpy \\
L6 & \freqmed{} & Coding & Implement LayerNorm forward and backward in numpy \\
L5 & \freqhigh{} & Coding & Write a minimal training loop with gradient clipping in PyTorch \\
L6 & \freqlow{} & Coding & Implement 2-D convolution as a matrix multiply (im2col) \\
L5 & \freqhigh{} & Coding & Implement BPE---train the merges on a corpus, then encode a word \\
L5 & \freqhigh{} & Coding & Implement sampling from logits with temperature, top-$k$, and top-$p$ \\
L6 & \freqhigh{} & Coding & Implement greedy decoding with a KV cache \\
L6 & \freqmed{} & Coding & Implement beam search with length normalization \\
L5 & \freqmed{} & Coding & Build a toy inverted index and score queries with BM25 \\
L5 & \freqhigh{} & Coding & Implement NDCG@k \\
L6 & \freqmed{} & Coding & Implement the in-batch softmax loss for a two-tower retrieval model \\
L6 & \freqmed{} & Coding & Implement greedy graph search over a fixed HNSW-style neighbor graph \\
L5 & \freqhigh{} & Coding & Implement logistic regression with SGD from scratch, deriving the gradient \\
L5 & \freqhigh{} & Coding & Implement k-means with k-means++ initialization \\
L6 & \freqmed{} & Coding & Implement gradient boosting with decision stumps \\
\bottomrule
\end{tabularx}

\subsection*{Chapter 29: The Staff Loop}
\begin{tabularx}{\textwidth}{@{}l l l X@{}}
\toprule
\textbf{Lvl} & \textbf{Freq} & \textbf{Type} & \textbf{Question} \\
\midrule
L6 & \freqhigh{} & Behavioral & Walk me through your most impactful project. (Then 40 minutes of drilling.) \\
L6 & \freqhigh{} & Behavioral & A peer team's launch is hurting your metric. Their leadership is celebrating it. What do you do? \\
L7 & \freqmed{} & Behavioral & If we gave you a year and a small team, what would you work on, and why? \\
L6 & \freqmed{} & Behavioral & Here's a paper claiming a new post-training method beats RLHF. Critique its evaluation \\
L6 & \freqmed{} & Behavioral & You're two days from a launch your team has crunched for, and an eval shows a small but real regression in a harm category. Walk me through what you actually do \\
L6 & \freqhigh{} & Behavioral & Describe committing to a technical decision you argued against. Would you do it again? \\
L7 & \freqmed{} & Behavioral & If you joined and we gave you no direction for your first month, what would you do---concretely, in week one? \\
L6 & \freqhigh{} & System Design & I'm going to give you a deliberately vague prompt---``design memory for our assistant.'' Run the first ten minutes \\
\bottomrule
\end{tabularx}

\section{By Round}
The rounds a loop is actually made of. Drill one column at a time.
\subsection*{Breadth (75 questions)}
\begin{itemize}[leftmargin=*]
\item \textbf{Ch.~1} \quad L5 \quad What are eigenvalues and why do they matter for training neural networks?
\item \textbf{Ch.~1} \quad L5 \quad Explain SVD and give three applications in machine learning
\item \textbf{Ch.~1} \quad L5 \quad Why does the chain rule matter for backpropagation? Walk through a 3-layer example
\item \textbf{Ch.~1} \quad L5 \quad What is KL divergence? Why is it not a true distance metric? When do you use it in practice?
\item \textbf{Ch.~1} \quad L6 \quad KL divergence is not symmetric. When does this matter in practice?
\item \textbf{Ch.~1} \quad L6 \quad Derive the gradient of softmax cross-entropy loss with respect to the logits
\item \textbf{Ch.~1} \quad L7 \quad What is the rank of the attention matrix and why does this matter for efficiency?
\item \textbf{Ch.~2} \quad L5 \quad Explain the bias-variance tradeoff. Does it apply to deep learning?
\item \textbf{Ch.~2} \quad L5 \quad Explain double descent. Why does test error decrease again in the over-parameterized regime?
\item \textbf{Ch.~2} \quad L7 \quad Neural scaling laws: what scales and what does not? When do they break?
\item \textbf{Ch.~2} \quad L6 \quad What is grokking and what does it tell us about generalization?
\item \textbf{Ch.~2} \quad L6 \quad Are emergent abilities in LLMs real, or an artifact of evaluation?
\item \textbf{Ch.~2} \quad L7 \quad Why do large language models generalize despite having more parameters than training examples?
\item \textbf{Ch.~3} \quad L5 \quad When do you use softmax cross-entropy vs.\ sigmoid BCE?
\item \textbf{Ch.~3} \quad L5 \quad Why does SimCLR need large batch sizes, and what are the alternatives?
\item \textbf{Ch.~4} \quad L5 \quad Explain the dying ReLU problem. How do modern activations solve it?
\item \textbf{Ch.~4} \quad L6 \quad Why don't Transformers use BatchNorm? Derive the reasoning from first principles
\item \textbf{Ch.~5} \quad L5 \quad Explain Flash Attention. Why is it faster even though it does the same computation?
\item \textbf{Ch.~5} \quad L5 \quad How does RoPE encode position information? Why is it preferred over learned positional embeddings?
\item \textbf{Ch.~5} \quad L5 \quad Why does attention scale by $1/\sqrt{d_k}$? What happens if you don't?
\item \textbf{Ch.~5} \quad L5 \quad Why is self-attention permutation equivariant? Why does this matter for Transformers?
\item \textbf{Ch.~5} \quad L5 \quad What changed between GPT-2 (2019) and LLaMA (2023) architecturally? Why each change?
\item \textbf{Ch.~6} \quad L6 \quad What is embedding collapse and how do you detect/prevent it?
\item \textbf{Ch.~6} \quad L5 \quad How does subword tokenization (BPE) interact with embedding quality?
\item \textbf{Ch.~7} \quad L5 \quad Cosine similarity vs dot product vs L2 distance for embeddings---when each?
\item \textbf{Ch.~8} \quad L5 \quad What makes a good augmentation policy for contrastive learning, and how does the choice of augmentations define the learned representation?
\item \textbf{Ch.~8} \quad L6 \quad Next-token prediction vs masked language modeling---why did GPT's autoregressive approach win for generative AI while BERT's bidirectional approach initially dominated NLU benchmarks?
\item \textbf{Ch.~8} \quad L5 \quad Why do non-contrastive methods like BYOL work without negative samples? Shouldn't they collapse to a trivial solution?
\item \textbf{Ch.~8} \quad L6 \quad Explain the InfoNCE loss and its connection to mutual information. Why does temperature matter?
\item \textbf{Ch.~9} \quad L7 \quad How does in-context learning work mechanistically? Why can large language models learn from examples in the prompt without any gradient updates?
\item \textbf{Ch.~9} \quad L6 \quad Why do decoder-only models dominate modern NLP despite encoder-decoder being theoretically more flexible?
\item \textbf{Ch.~9} \quad L6 \quad Walk me through the DPO loss function. How does it eliminate the need for a separate reward model?
\item \textbf{Ch.~9} \quad L6 \quad Compare RoPE vs learned positional embeddings vs sinusoidal encodings. Why did RoPE become the dominant choice for modern LLMs?
\item \textbf{Ch.~10} \quad L5 \quad What is the effective receptive field of a CNN and why does it matter for architecture design?
\item \textbf{Ch.~10} \quad L5 \quad Explain how DETR works and why it represents a paradigm shift in object detection
\item \textbf{Ch.~11} \quad L5 \quad Explain the reparameterization trick and why it's needed for VAE training. What happens without it?
\item \textbf{Ch.~11} \quad L5 \quad What is classifier-free guidance and why does it improve generation quality in diffusion models?
\item \textbf{Ch.~11} \quad L5 \quad Explain the ELBO derivation for VAEs and why the model produces blurry images. What fixes exist?
\item \textbf{Ch.~11} \quad L7 \quad Flow matching vs denoising diffusion---what changed and why is it considered the next evolution?
\item \textbf{Ch.~11} \quad L5 \quad How does Stable Diffusion work end-to-end? Walk through the architecture from text prompt to generated image
\item \textbf{Ch.~11} \quad L6 \quad Explain the WGAN loss and why Wasserstein distance solves GAN training instability
\item \textbf{Ch.~12} \quad L5 \quad How do you handle the cold start problem for new users and new items?
\item \textbf{Ch.~13} \quad L6 \quad What are the theoretical expressivity limitations of message-passing GNNs, and when do they matter in practice?
\item \textbf{Ch.~13} \quad L5 \quad How do you handle heterogeneous graphs (multiple node types, multiple edge types) in a GNN?
\item \textbf{Ch.~15} \quad L5 \quad Explain gradient accumulation and when you would use it
\item \textbf{Ch.~16} \quad L5 \quad What is negative transfer and how do you detect and prevent it?
\item \textbf{Ch.~16} \quad L6 \quad Why does LoRA work? What is the low-rank assumption about fine-tuning?
\item \textbf{Ch.~17} \quad L5 \quad Walk me through the RLHF pipeline for aligning an LLM
\item \textbf{Ch.~17} \quad L5 \quad What is reward hacking? Give 3 concrete examples and how to mitigate each
\item \textbf{Ch.~17} \quad L5 \quad Why does the KL penalty matter in RLHF? What happens without it?
\item \textbf{Ch.~17} \quad L6 \quad Explain GRPO. Why did it replace PPO's critic for reasoning RL, and what breaks when you use it naively?
\item \textbf{Ch.~18} \quad L5 \quad How does CLIP's contrastive loss create a shared embedding space?
\item \textbf{Ch.~18} \quad L5 \quad CLIP vs SigLIP---what changed and why?
\item \textbf{Ch.~18} \quad L5 \quad Compare early, late, and cross-attention fusion for multimodal models. When would you use each?
\item \textbf{Ch.~19} \quad L5 \quad Explain speculative decoding. When would you use it and when would you not?
\item \textbf{Ch.~20} \quad L5 \quad Jailbreak attacks: how do they work and how do you defend?
\item \textbf{Ch.~20} \quad L6 \quad What is mechanistic interpretability? Give a concrete example of a discovered circuit
\item \textbf{Ch.~20} \quad L5 \quad How do you detect and reduce hallucination in a production LLM system?
\item \textbf{Ch.~20} \quad L7 \quad How can you use representation engineering (steering vectors) to control model behavior without retraining?
\item \textbf{Ch.~22} \quad L5 \quad How do you handle training-serving skew?
\item \textbf{Ch.~22} \quad L5 \quad How do you detect and handle concept drift vs data drift?
\item \textbf{Ch.~23} \quad L5 \quad Your model has 95\% accuracy but stakeholders are unhappy. What is going on?
\item \textbf{Ch.~23} \quad L5 \quad How do you evaluate a search ranking system?
\item \textbf{Ch.~23} \quad L5 \quad How do you know if a model improvement is statistically significant?
\item \textbf{Ch.~23} \quad L5 \quad What is wrong with accuracy as a metric? When is it actually fine?
\item \textbf{Ch.~25} \quad L5 \quad Walk me through what happens between ``we have a Common Crawl snapshot'' and the first training step
\item \textbf{Ch.~25} \quad L5 \quad Why does deduplication improve language models, and what goes wrong if you deduplicate too aggressively?
\item \textbf{Ch.~25} \quad L6 \quad What is mid-training (annealing), why does everyone do it now, and what is the catch?
\item \textbf{Ch.~26} \quad L5 \quad Walk me through the GPU memory hierarchy. Why does it, rather than FLOPs, dominate ML performance thinking?
\item \textbf{Ch.~26} \quad L5 \quad Why do tensor cores only accelerate matrix multiplication? What does that imply for how you design models and kernels?
\item \textbf{Ch.~26} \quad L6 \quad Kernel fusion doesn't reduce FLOPs. Why does it make models faster---and when does it stop helping? Use FlashAttention as your example
\item \textbf{Ch.~27} \quad L5 \quad What makes an LLM system an ``agent''? Walk me through the agent loop and where the engineering difficulty actually lives
\item \textbf{Ch.~27} \quad L5 \quad Explain pass@$k$ vs.\ pass\textasciicircum{}$k$. Your agent team reports 92\% pass@5---what do you tell the VP who wants to ship?
\item \textbf{Ch.~27} \quad L6 \quad Why do computer-use agents lag API-based agents so badly, and when would you deploy one anyway?
\item \textbf{Ch.~27} \quad L5 \quad What does MCP actually standardize, and what does it change---and not change---about how you architect an agent?
\end{itemize}

\subsection*{Depth (153 questions)}
\begin{itemize}[leftmargin=*]
\item \textbf{Ch.~1} \quad L5 \quad Estimate the FLOPs for multiplying two matrices of size $(1024, 4096)$ and $(4096, 1024)$
\item \textbf{Ch.~1} \quad L6 \quad Your gradient norms are 1000$\times$ larger in early layers than later layers. What is happening mathematically?
\item \textbf{Ch.~1} \quad L5 \quad Estimate the memory for storing a 50K $\times$ 768 embedding matrix in FP32, FP16, and INT8
\item \textbf{Ch.~2} \quad L6 \quad Your 10B parameter model is clearly overparameterized for your task. Should you reduce it?
\item \textbf{Ch.~2} \quad L6 \quad Your team has a fixed compute budget. How do you decide model size vs. training data?
\item \textbf{Ch.~2} \quad L6 \quad Estimate the compute-optimal model size for a training budget of $10^{22}$ FLOPs using Chinchilla scaling laws
\item \textbf{Ch.~2} \quad L6 \quad Training loss is still decreasing but validation loss plateaued. You have $10\times$ more unlabeled data available. What do you do?
\item \textbf{Ch.~3} \quad L6 \quad Your model's training loss is decreasing but validation loss curves show the model is increasingly miscalibrated. What loss function issue could cause this?
\item \textbf{Ch.~3} \quad L5 \quad Cross-entropy vs.\ focal loss vs.\ class-balanced loss for imbalanced classification---what is your decision framework?
\item \textbf{Ch.~3} \quad L6 \quad Your SFT loss changes when you change batch composition or gradient-accumulation steps---same data, same model, same effective batch size. Why?
\item \textbf{Ch.~3} \quad L6 \quad Estimate the memory and FLOP cost of the cross-entropy layer for a 128K-vocabulary model training on a 1M-token batch. How do you reduce it?
\item \textbf{Ch.~3} \quad L6 \quad Your contrastive learning model's loss suddenly jumps mid-training and doesn't recover. Diagnose
\item \textbf{Ch.~3} \quad L6 \quad Pointwise vs.\ pairwise vs.\ listwise ranking losses---when do you use each?
\item \textbf{Ch.~3} \quad L5 \quad Your classification model gets 95\% accuracy but the loss curve shows the model is still decreasing slowly. Should you keep training?
\item \textbf{Ch.~4} \quad L5 \quad Why GELU over ReLU? Why SwiGLU over GELU?
\item \textbf{Ch.~4} \quad L6 \quad Your model shows different behavior at training vs inference. What normalization issue could cause this?
\item \textbf{Ch.~4} \quad L5 \quad RMSNorm vs LayerNorm---when would you choose each?
\item \textbf{Ch.~4} \quad L6 \quad Pre-norm vs. post-norm---derive the gradient-flow difference and explain what modern hybrids fix
\item \textbf{Ch.~4} \quad L5 \quad Training loss is NaN after a few hundred steps. Which activation/normalization issues could cause this?
\item \textbf{Ch.~4} \quad L6 \quad Your 70B-parameter run in BF16 starts loss-spiking at 100B tokens. Walk me through the activation- and normalization-level causes and mitigations
\item \textbf{Ch.~5} \quad L6 \quad Compare MHA, GQA, MQA, and MLA for KV-cache efficiency. Why does MLA need a decoupled RoPE key?
\item \textbf{Ch.~5} \quad L6 \quad Estimate the memory required to train a 7B parameter model with batch size 1, sequence length 2048
\item \textbf{Ch.~5} \quad L5 \quad Compare MHA, GQA, and MQA. When would you choose each?
\item \textbf{Ch.~5} \quad L6 \quad Your large BF16 pretraining run shows intermittent loss spikes. Instrumentation shows attention logits growing into the hundreds over training. Diagnose and fix
\item \textbf{Ch.~5} \quad L6 \quad Calculate the FLOPs for a single forward pass through a 7B parameter Transformer
\item \textbf{Ch.~5} \quad L6 \quad Calculate the KV cache memory for serving LLaMA 70B at 128K context length
\item \textbf{Ch.~5} \quad L5 \quad Pre-norm vs Post-norm Transformers --- which would you choose and why?
\item \textbf{Ch.~5} \quad L6 \quad Your Transformer model shows degrading perplexity at 32K context despite being trained on 8K. Diagnose and fix
\item \textbf{Ch.~5} \quad L6 \quad Compare Transformer vs SSM (Mamba) architectures. When would you choose each?
\item \textbf{Ch.~5} \quad L7 \quad Your model's attention patterns show all heads attending to the same positions (typically position 0 or EOS). What's happening?
\item \textbf{Ch.~5} \quad L7 \quad Estimate the ratio of attention FLOPs to FFN FLOPs at different sequence lengths. When does attention become the bottleneck?
\item \textbf{Ch.~5} \quad L6 \quad Flash Attention vs Linear Attention --- when would you use each?
\item \textbf{Ch.~6} \quad L6 \quad Your embedding space shows that unrelated items cluster together. Diagnose and fix
\item \textbf{Ch.~6} \quad L6 \quad You are putting 100M item embeddings behind a vector index. Which embedding-side choices do you make---dimension, precision, refresh---and how would you know if one of them silently cost you recall?
\item \textbf{Ch.~6} \quad L5 \quad Estimate the memory for an embedding table with 50K vocab, 768 dimensions in FP16 vs INT8
\item \textbf{Ch.~6} \quad L5 \quad How would you debug poor embedding quality?
\item \textbf{Ch.~7} \quad L5 \quad Your two-tower model retrieves documents that share keywords with the query but aren't actually relevant. How do you fix this?
\item \textbf{Ch.~7} \quad L5 \quad Siamese vs Triplet vs Two-Tower---decision framework for a new retrieval project
\item \textbf{Ch.~7} \quad L6 \quad Your two-tower model's recall@10 is good but precision@10 is poor. What's wrong?
\item \textbf{Ch.~7} \quad L5 \quad Cross-encoder vs bi-encoder---when is the quality gap worth the latency cost?
\item \textbf{Ch.~7} \quad L6 \quad Hard negative mining strategies---how to choose and when each fails
\item \textbf{Ch.~7} \quad L7 \quad Your contrastive model produces good embeddings for frequent items but poor for rare items. Why?
\item \textbf{Ch.~7} \quad L6 \quad Estimate the QPS of a retrieval system using HNSW on 100M documents. Then add a cross-encoder reranker over the top 100---what changes?
\item \textbf{Ch.~8} \quad L5 \quad SimCLR vs BYOL vs MAE---give me a decision framework for choosing between them
\item \textbf{Ch.~8} \quad L6 \quad Your SSL model's downstream performance plateaus despite increasing pre-training data from 1M to 10M images. What could be going wrong?
\item \textbf{Ch.~8} \quad L5 \quad BERT vs GPT pre-training objectives---what are the fundamental tradeoffs between masked language modeling and autoregressive language modeling?
\item \textbf{Ch.~8} \quad L6 \quad Estimate the compute cost of pre-training a BERT-base model on 16GB of text data
\item \textbf{Ch.~9} \quad L5 \quad You're told a model is ``Llama-class, 7B parameters.'' Roughly where do those 7B parameters live?
\item \textbf{Ch.~9} \quad L5 \quad Encoder-only vs decoder-only vs encoder-decoder---when do you use each, and what drives the decision?
\item \textbf{Ch.~9} \quad L6 \quad Your fine-tuned chat model performs worse than the base model it was tuned from. Walk me through your diagnosis
\item \textbf{Ch.~9} \quad L6 \quad Your fine-tuned LLM generates fluent but factually wrong answers. Diagnose the issue and propose solutions
\item \textbf{Ch.~9} \quad L6 \quad Estimate the inference latency for generating 100 tokens with a 7B parameter model on a single A100 GPU
\item \textbf{Ch.~9} \quad L5 \quad BPE vs SentencePiece vs character-level tokenization---what are the tradeoffs and how does tokenizer choice affect model performance?
\item \textbf{Ch.~10} \quad L5 \quad CNN vs ViT for a dataset with only 10K labeled images---how do you decide, and what changes with 10M images?
\item \textbf{Ch.~10} \quad L6 \quad Estimate the FLOPs for a ResNet-50 forward pass on a $224 \times 224$ image. How does this compare to ViT-B/16?
\item \textbf{Ch.~10} \quad L6 \quad Your ViT model achieves 92\% accuracy on ImageNet but only 65\% on your domain-specific dataset despite fine-tuning. Diagnose the problem
\item \textbf{Ch.~10} \quad L6 \quad Compare the ResNet skip connection to DenseNet dense connections to Swin Transformer's shifted windows. What problem does each solve?
\item \textbf{Ch.~11} \quad L6 \quad Diffusion models vs GANs vs VAEs vs autoregressive models---give me a comparison table and decision framework for different generation tasks
\item \textbf{Ch.~11} \quad L6 \quad Your diffusion model generates high-quality images overall but has mode collapse for certain categories (e.g., generates dogs well but all cats look the same). Diagnose and fix
\item \textbf{Ch.~11} \quad L7 \quad Estimate the inference cost (FLOPs and latency) for generating a $512 \times 512$ image with Stable Diffusion on an A100 GPU
\item \textbf{Ch.~12} \quad L5 \quad Two-tower vs.\ interaction-based models for candidate generation---when would you use each?
\item \textbf{Ch.~13} \quad L6 \quad When should you NOT use a GNN? Give examples where simpler approaches win
\item \textbf{Ch.~13} \quad L5 \quad Your GNN's performance degrades as you add more layers. What's happening and how do you fix it?
\item \textbf{Ch.~13} \quad L5 \quad GCN vs GraphSAGE vs GAT---when would you use each?
\item \textbf{Ch.~13} \quad L6 \quad Estimate the memory requirements for running a 3-layer GCN on a graph with 10M nodes
\item \textbf{Ch.~13} \quad L6 \quad GCN vs GraphSAGE vs GAT---which would you choose for a production recommendation system at scale?
\item \textbf{Ch.~14} \quad L6 \quad When would you choose MoE over dense scaling, and vice versa?
\item \textbf{Ch.~14} \quad L6 \quad You're building a system that processes 100K-token documents. Compare transformer vs Mamba for this use case
\item \textbf{Ch.~14} \quad L6 \quad Estimate the compute savings of a MoE model with 8 experts, top-2 routing, vs a dense model of the same quality
\item \textbf{Ch.~14} \quad L6 \quad INT8 vs INT4 vs FP8 quantization---what is your decision framework?
\item \textbf{Ch.~14} \quad L6 \quad Your quantized model performs well on benchmarks but fails on specific domains (code, math). Why?
\item \textbf{Ch.~14} \quad L6 \quad MoE routing collapse: what is it, how do you detect it, and how do you fix it?
\item \textbf{Ch.~14} \quad L6 \quad Structured vs unstructured pruning---which actually gives real speedup?
\item \textbf{Ch.~15} \quad L5 \quad Your model's loss plateaus after initial descent. Walk through your debugging process
\item \textbf{Ch.~15} \quad L5 \quad Adam vs SGD with momentum---when would you pick each?
\item \textbf{Ch.~15} \quad L6 \quad Your model's loss spikes every N steps then recovers. Diagnose
\item \textbf{Ch.~15} \quad L5 \quad Training loss is NaN after 1000 steps. Walk through your debugging process
\item \textbf{Ch.~15} \quad L7 \quad Your 70B pretraining run diverges at 200B tokens. Give a full differential diagnosis, ordered by prior probability, with the cheapest discriminating experiment for each
\item \textbf{Ch.~15} \quad L6 \quad Your distributed training job is 60\% as fast as expected with 8 GPUs. What is wrong?
\item \textbf{Ch.~15} \quad L7 \quad A node dies every 3 hours on your 16K-GPU run. Design the recovery system
\item \textbf{Ch.~15} \quad L6 \quad How much GPU memory is needed to fine-tune a 13B parameter model with LoRA vs.\ full fine-tuning?
\item \textbf{Ch.~15} \quad L7 \quad Calculate the training throughput (tokens/sec) for a 7B model on 8$\times$A100s
\item \textbf{Ch.~15} \quad L5 \quad Cosine decay vs.\ constant LR with warmup---when would you use each?
\item \textbf{Ch.~15} \quad L6 \quad Data parallelism vs.\ tensor parallelism vs.\ pipeline parallelism---what is your decision framework?
\item \textbf{Ch.~15} \quad L6 \quad FP16 vs.\ BF16 vs.\ FP8 for training---when would you use each?
\item \textbf{Ch.~16} \quad L5 \quad Estimate the number of trainable parameters for LoRA with rank 16 on a 7B model
\item \textbf{Ch.~16} \quad L6 \quad Your LoRA fine-tuned model performs well on your test set but catastrophically forgets general capabilities. How do you fix it?
\item \textbf{Ch.~16} \quad L6 \quad Your fine-tuned model echoes the prompt before answering---why?
\item \textbf{Ch.~16} \quad L6 \quad LoRA vs QLoRA vs full fine-tuning vs prompt tuning---give a complete decision framework
\item \textbf{Ch.~16} \quad L6 \quad You have chosen LoRA to adapt a 7B instruction-tuned model to an internal support-and-policy domain: 20K curated examples, one 80\,GB GPU, and the adapter will be served alongside the general assistant. Pick rank, $\alpha$, and target modules, and defend the choices
\item \textbf{Ch.~16} \quad L5 \quad Adapter layers vs LoRA vs prefix tuning---what are the architectural differences and when would you use each?
\item \textbf{Ch.~16} \quad L6 \quad You fine-tuned a model and it performs well on your test set but poorly in production. What happened?
\item \textbf{Ch.~17} \quad L6 \quad PPO vs DPO vs RLHF---complete comparison for LLM alignment
\item \textbf{Ch.~17} \quad L6 \quad Your RLHF-tuned model becomes sycophantic (agrees with everything). Diagnose and fix
\item \textbf{Ch.~17} \quad L6 \quad Estimate the memory required to run RLHF (4 models: policy, reference, reward model, critic)
\item \textbf{Ch.~17} \quad L6 \quad Constitutional AI vs RLHF vs DPO---when would you use each?
\item \textbf{Ch.~17} \quad L7 \quad Your RLHF training shows the reward score plateauing while KL divergence keeps climbing. What is happening?
\item \textbf{Ch.~17} \quad L6 \quad Process reward models vs outcome reward models---what are the tradeoffs and when does each win?
\item \textbf{Ch.~17} \quad L6 \quad Your GRPO run's mean reward climbs but responses grow unboundedly long and benchmark scores stall---diagnose
\item \textbf{Ch.~17} \quad L7 \quad When is test-time compute cheaper than a bigger model---and when does it stop working?
\item \textbf{Ch.~18} \quad L6 \quad Cascade vs native speech-to-speech for a voice assistant---decide and defend
\item \textbf{Ch.~18} \quad L6 \quad Adapter-based (LLaVA) vs natively multimodal (Gemini)---tradeoffs
\item \textbf{Ch.~18} \quad L7 \quad Estimate the compute cost of training a CLIP model on 400M image-text pairs
\item \textbf{Ch.~18} \quad L6 \quad Your VLM hallucinates visual details that are not in the image. Why and how to fix?
\item \textbf{Ch.~18} \quad L6 \quad Visual tokens consume 576 tokens of context. How do you reduce this cost?
\item \textbf{Ch.~19} \quad L6 \quad KV cache memory is your bottleneck---you are running out of GPU memory and dropping requests. What do you do?
\item \textbf{Ch.~19} \quad L7 \quad Calculate the throughput (tokens/sec) of a 70B model on 8$\times$H100s with tensor parallelism
\item \textbf{Ch.~19} \quad L6 \quad How much memory savings does INT4 quantization give for a 70B model? What is the quality trade-off?
\item \textbf{Ch.~19} \quad L7 \quad Your LLM service needs to handle 500 concurrent users with 2-second time-to-first-token. Size the infrastructure
\item \textbf{Ch.~19} \quad L6 \quad GPTQ vs AWQ vs GGUF quantization---when would you use each?
\item \textbf{Ch.~19} \quad L6 \quad Tensor parallelism vs pipeline parallelism for inference---what is your decision framework?
\item \textbf{Ch.~19} \quad L7 \quad Prefill-optimized vs decode-optimized serving---when should you disaggregate them?
\item \textbf{Ch.~19} \quad L6 \quad Your LLM serving system's p99 latency is 10$\times$ the p50. Diagnose
\item \textbf{Ch.~19} \quad L6 \quad After quantizing your model to INT4, certain types of prompts give garbage output while others are fine. Why?
\item \textbf{Ch.~19} \quad L6 \quad Your speculative decoding setup is slower than standard decoding. What went wrong?
\item \textbf{Ch.~20} \quad L7 \quad Your LLM refuses too many benign queries. How do you reduce over-refusal without increasing safety risk?
\item \textbf{Ch.~20} \quad L6 \quad A jailbreak goes viral against your production model---walk me through the first 24 hours
\item \textbf{Ch.~21} \quad L6 \quad Your 8K-trained model must serve 64K contexts next month. What are your options, and what do they cost?
\item \textbf{Ch.~21} \quad L6 \quad Your model's quality degrades beyond 32K context. Walk me through the diagnosis
\item \textbf{Ch.~21} \quad L6 \quad Estimate the cost per query for a RAG system vs a long-context LLM (128K)
\item \textbf{Ch.~21} \quad L5 \quad Long context window vs RAG---when do you choose each?
\item \textbf{Ch.~22} \quad L6 \quad Your model's CTR dropped 5\% overnight. Walk through your diagnosis
\item \textbf{Ch.~22} \quad L6 \quad Your model's online metrics diverge from offline metrics after 2 weeks in production. What is happening?
\item \textbf{Ch.~22} \quad L7 \quad Estimate the infrastructure cost for serving a recommendation model at 100K QPS
\item \textbf{Ch.~22} \quad L5 \quad Shadow deployment vs A/B test vs multi-armed bandit---when do you use each?
\item \textbf{Ch.~23} \quad L6 \quad Model B is 1.5 points better than model A on MMLU---do you ship it?
\item \textbf{Ch.~23} \quad L5 \quad Precision is 0.95 but recall is 0.30. What do you do?
\item \textbf{Ch.~23} \quad L6 \quad Your model's AUC is 0.95 but calibration is poor (ECE = 0.15). Why does this matter and how do you fix it?
\item \textbf{Ch.~23} \quad L5 \quad NDCG vs MAP vs MRR for ranking---when do you use each?
\item \textbf{Ch.~23} \quad L6 \quad Estimate the sample size needed for an A/B test to detect a 1\% improvement in CTR
\item \textbf{Ch.~23} \quad L6 \quad Offline metrics went up but online metrics went down. Give 5 possible reasons
\item \textbf{Ch.~24} \quad L6 \quad Your model needs to run in $<$10ms. Current latency is 50ms. Walk through the optimization decision tree
\item \textbf{Ch.~24} \quad L6 \quad You have 3 months and a team of 3 ML engineers. Should you fine-tune an LLM or train a custom model?
\item \textbf{Ch.~24} \quad L5 \quad GPU vs TPU vs CPU for inference---give a decision framework for different workloads
\item \textbf{Ch.~24} \quad L7 \quad Build vs buy vs open-source---give a decision framework for ML infrastructure
\item \textbf{Ch.~24} \quad L5 \quad Would you use BERT or GPT for sentiment analysis?
\item \textbf{Ch.~24} \quad L5 \quad When should you use gradient boosting vs.\ neural networks for tabular data?
\item \textbf{Ch.~25} \quad L6 \quad How would you decide the code:web:math mixture for a pretraining run?
\item \textbf{Ch.~25} \quad L6 \quad After a data-pipeline change, training loss improved---but downstream evals got worse. Diagnose
\item \textbf{Ch.~25} \quad L7 \quad Your data lead proposes raising the synthetic fraction of the next run from 10\% to 40\%. How do you evaluate the proposal?
\item \textbf{Ch.~25} \quad L6 \quad You are choosing between a 32K, 128K, and 256K vocabulary for a new model family. What changes, and when would you retrain vs.\ reuse a tokenizer?
\item \textbf{Ch.~25} \quad L6 \quad At 40\% through a flagship run, held-out code loss jumps and stays elevated. Walk me through the response
\item \textbf{Ch.~26} \quad L6 \quad I give you an op---say, a LayerNorm over a $[16{,}384 \times 8{,}192]$ BF16 tensor, or a $4096^3$ GEMM. Is it compute-bound or bandwidth-bound on an H100? Show your method
\item \textbf{Ch.~26} \quad L6 \quad Your training job shows the GPU at 20\% utilization. Diagnose it
\item \textbf{Ch.~26} \quad L6 \quad How long does it take to all-reduce the gradients of a 70B-parameter model across 8 GPUs? Derive it
\item \textbf{Ch.~26} \quad L6 \quad A teammate changed the hidden size from 4096 to 4100 ``to add a few features,'' and training throughput dropped 35\%. What happened?
\item \textbf{Ch.~26} \quad L6 \quad What problem do CUDA graphs solve, and what constraints do they impose? Where do they matter most in an LLM system?
\item \textbf{Ch.~26} \quad L7 \quad Estimate the maximum tokens/second for a 70B dense model on 8$\times$H100---for training, and for inference
\item \textbf{Ch.~26} \quad L7 \quad Your 512-GPU pretraining run reports 25\% MFU. Walk me through how you find the missing performance
\item \textbf{Ch.~27} \quad L6 \quad When should you NOT build an agent? Your PM wants ``an agent'' for a document-processing product---how do you decide?
\item \textbf{Ch.~27} \quad L7 \quad Your agent succeeds on 60\% of an internal task suite. Leadership wants 95\%. Take it there
\item \textbf{Ch.~27} \quad L6 \quad One agent with 30 tools, or an orchestrator with specialized sub-agents? Argue both sides with numbers
\item \textbf{Ch.~27} \quad L6 \quad Your agent product works but loses money: \$4 of inference per task against \$1.50 of revenue. Fix the economics without wrecking the success rate
\end{itemize}

\subsection*{Coding (16 questions)}
\begin{itemize}[leftmargin=*]
\item \textbf{Ch.~28} \quad L5 \quad Implement multi-head self-attention with a causal mask in numpy
\item \textbf{Ch.~28} \quad L5 \quad Implement one AdamW update step in numpy
\item \textbf{Ch.~28} \quad L6 \quad Implement LayerNorm forward and backward in numpy
\item \textbf{Ch.~28} \quad L5 \quad Write a minimal training loop with gradient clipping in PyTorch
\item \textbf{Ch.~28} \quad L6 \quad Implement 2-D convolution as a matrix multiply (im2col)
\item \textbf{Ch.~28} \quad L5 \quad Implement BPE---train the merges on a corpus, then encode a word
\item \textbf{Ch.~28} \quad L5 \quad Implement sampling from logits with temperature, top-$k$, and top-$p$
\item \textbf{Ch.~28} \quad L6 \quad Implement greedy decoding with a KV cache
\item \textbf{Ch.~28} \quad L6 \quad Implement beam search with length normalization
\item \textbf{Ch.~28} \quad L5 \quad Build a toy inverted index and score queries with BM25
\item \textbf{Ch.~28} \quad L5 \quad Implement NDCG@k
\item \textbf{Ch.~28} \quad L6 \quad Implement the in-batch softmax loss for a two-tower retrieval model
\item \textbf{Ch.~28} \quad L6 \quad Implement greedy graph search over a fixed HNSW-style neighbor graph
\item \textbf{Ch.~28} \quad L5 \quad Implement logistic regression with SGD from scratch, deriving the gradient
\item \textbf{Ch.~28} \quad L5 \quad Implement k-means with k-means++ initialization
\item \textbf{Ch.~28} \quad L6 \quad Implement gradient boosting with decision stumps
\end{itemize}

\subsection*{Design (50 questions)}
\begin{itemize}[leftmargin=*]
\item \textbf{Ch.~2} \quad L7 \quad You get 1\% of the target training budget to run scaling experiments that must predict the final run's loss. Design the study
\item \textbf{Ch.~3} \quad L6 \quad You're building a visual search system for an e-commerce platform with 10M products. What loss function would you use for training the embedding model?
\item \textbf{Ch.~3} \quad L5 \quad You need to train an embedding model for semantic search across 100M documents. Which loss function and why?
\item \textbf{Ch.~3} \quad L6 \quad Design the loss function for a retrieval system where you have click data (implicit feedback) but no explicit relevance labels
\item \textbf{Ch.~3} \quad L5 \quad You're building a product recommendation system with implicit feedback (clicks). What loss would you use and why?
\item \textbf{Ch.~3} \quad L7 \quad You're building a multi-task model that does classification, regression, and ranking simultaneously. Design the loss function
\item \textbf{Ch.~5} \quad L7 \quad Design the attention mechanism for a model that must process 1 million token documents
\item \textbf{Ch.~5} \quad L7 \quad You need to modify a standard Transformer for real-time streaming applications (e.g., live transcription). What changes?
\item \textbf{Ch.~8} \quad L6 \quad Design an SSL pre-training pipeline for a dataset of 10M unlabeled images and 100K labeled images. Walk through method selection, architecture, training, and evaluation
\item \textbf{Ch.~8} \quad L6 \quad You have 10M unlabeled images and 10K labeled images in a medical imaging domain. How would you use SSL to improve your classifier, and what challenges do you expect?
\item \textbf{Ch.~9} \quad L7 \quad Design the architecture for a multilingual document understanding system that handles 50+ languages, mixed-language documents, and both structured (tables, forms) and unstructured text
\item \textbf{Ch.~9} \quad L6 \quad Design the architecture for a customer support ticket classification system with 50 categories, handling multilingual input and evolving category taxonomy
\item \textbf{Ch.~10} \quad L6 \quad Design a vision backbone for real-time object detection on an edge device (e.g., mobile phone or drone with $<$10W power budget)
\item \textbf{Ch.~10} \quad L6 \quad Design a visual search system that handles 100M product images with sub-200ms latency. Walk through the architecture
\item \textbf{Ch.~11} \quad L6 \quad Design an image generation pipeline for a product photography application (e.g., placing products in realistic scenes)
\item \textbf{Ch.~13} \quad L6 \quad Design a GNN-based system for real-time fraud detection on a financial transaction graph
\item \textbf{Ch.~14} \quad L6 \quad Design an efficient architecture for on-device ML (phone, $<$100ms latency, $<$50MB model)
\item \textbf{Ch.~14} \quad L7 \quad You have a serving budget of 80\,GB of GPU memory and 20\,ms/token decode latency, with contexts up to 128K. Design the model
\item \textbf{Ch.~14} \quad L6 \quad Knowledge distillation: design a pipeline to compress a 70B teacher to a 7B student
\item \textbf{Ch.~14} \quad L7 \quad Design an efficient architecture for real-time video understanding on edge devices
\item \textbf{Ch.~15} \quad L6 \quad How would you scale training of a 70B parameter model across 256 GPUs?
\item \textbf{Ch.~16} \quad L6 \quad Design a fine-tuning pipeline for adapting a foundation model to a new domain with 10K examples
\item \textbf{Ch.~16} \quad L5 \quad You need to fine-tune a 70B LLM for a customer's specific domain. You have access to a single A100 80GB GPU. What approach would you take?
\item \textbf{Ch.~16} \quad L6 \quad You have one base model and 100 customers, each with their own LoRA adapter. Design the serving system
\item \textbf{Ch.~17} \quad L6 \quad Design a reward model training pipeline. What data do you need?
\item \textbf{Ch.~18} \quad L6 \quad Design a multimodal content understanding system for a social media platform
\item \textbf{Ch.~18} \quad L6 \quad Design a multimodal search system that lets users search a product catalog using text or image queries
\item \textbf{Ch.~18} \quad L6 \quad How would you add audio understanding to an existing vision-language model?
\item \textbf{Ch.~19} \quad L7 \quad Design an LLM serving stack that handles 1,000 queries per second with a p99 latency of 2 seconds for a 70B parameter model
\item \textbf{Ch.~19} \quad L7 \quad Your capacity plan assumed 500-token outputs. The reasoning model you are now deploying emits $\sim$20K thinking tokens per query. Redo the plan
\item \textbf{Ch.~20} \quad L6 \quad Design a content filtering pipeline for an LLM chatbot with $<$50ms added latency
\item \textbf{Ch.~20} \quad L6 \quad Hallucination detection: design a system that flags when an LLM generates unsupported claims
\item \textbf{Ch.~20} \quad L6 \quad How do you set up a red teaming program for an LLM product?
\item \textbf{Ch.~20} \quad L7 \quad Design safety for an agent that can browse the web and execute code
\item \textbf{Ch.~22} \quad L6 \quad Design ML infrastructure for a search ranking system at 10K QPS
\item \textbf{Ch.~22} \quad L6 \quad How would you set up an A/B test for a new recommendation model?
\item \textbf{Ch.~22} \quad L6 \quad Design the monitoring system for an ML model in production
\item \textbf{Ch.~22} \quad L6 \quad Feature store: batch vs streaming features---describe the architecture and tradeoffs
\item \textbf{Ch.~23} \quad L7 \quad Design the eval stack that decides whether tomorrow's checkpoint ships
\item \textbf{Ch.~23} \quad L7 \quad Design an evaluation framework for an LLM-based chatbot
\item \textbf{Ch.~24} \quad L6 \quad Design a content recommendation system for a news app
\item \textbf{Ch.~24} \quad L6 \quad Design the ML architecture for a content moderation system
\item \textbf{Ch.~25} \quad L7 \quad You own data for the next frontier pretraining run. Design the pipeline and the ablation program
\item \textbf{Ch.~25} \quad L6 \quad Design the benchmark decontamination strategy for a pretraining run. What does n-gram overlap miss, and what do you do about it?
\item \textbf{Ch.~25} \quad L6 \quad Design the in-flight evaluation for a three-month pretraining run. When do in-loop evals mislead?
\item \textbf{Ch.~26} \quad L7 \quad You have 8 nodes of 8$\times$H100. Map tensor, pipeline, data, and expert parallelism onto this cluster's fabric, justifying the mapping from link characteristics---then tell me what changes on a GB200 NVL72
\item \textbf{Ch.~27} \quad L6 \quad Design a code-review agent for your organization's monorepo
\item \textbf{Ch.~27} \quad L6 \quad Design the evaluation for a customer-support agent that can read accounts, process refunds, and escalate to humans
\item \textbf{Ch.~27} \quad L7 \quad Your agent's tasks run for hours and blow past the context window. Design its memory and state management
\item \textbf{Ch.~29} \quad L6 \quad I'm going to give you a deliberately vague prompt---``design memory for our assistant.'' Run the first ten minutes
\end{itemize}

\subsection*{Behavioral (7 questions)}
\begin{itemize}[leftmargin=*]
\item \textbf{Ch.~29} \quad L6 \quad Walk me through your most impactful project. (Then 40 minutes of drilling.)
\item \textbf{Ch.~29} \quad L6 \quad A peer team's launch is hurting your metric. Their leadership is celebrating it. What do you do?
\item \textbf{Ch.~29} \quad L7 \quad If we gave you a year and a small team, what would you work on, and why?
\item \textbf{Ch.~29} \quad L6 \quad Here's a paper claiming a new post-training method beats RLHF. Critique its evaluation
\item \textbf{Ch.~29} \quad L6 \quad You're two days from a launch your team has crunched for, and an eval shows a small but real regression in a harm category. Walk me through what you actually do
\item \textbf{Ch.~29} \quad L6 \quad Describe committing to a technical decision you argued against. Would you do it again?
\item \textbf{Ch.~29} \quad L7 \quad If you joined and we gave you no direction for your first month, what would you do---concretely, in week one?
\end{itemize}

